% Think time: the paper. The prose follows docs/writing.md. Every number,
% table, figure and board is written by `lab paper think-time` into
% generated/, from the model's committed results, and read back with \V{key}
% and \block{name}.
\documentclass[11pt]{article}
% The preamble every paper of the lab shares: % The preamble every paper of the lab shares: % The preamble every paper of the lab shares: \input{../../../papers/preamble}
% from models/<model>/paper/, \input{../preamble} from papers/<name>/.
% Boards come from chessboard and xskak, plots from pgfplots. A paper prints no
% number by hand: `lab paper` writes every number, table and figure under the
% paper's generated/ folder, and paper.tex says what they mean.
\usepackage[LSBC4,T1]{fontenc}
\usepackage[utf8]{inputenc}
\usepackage{lmodern}
\usepackage[margin=1in]{geometry}
\usepackage{microtype}
\usepackage{amsmath}
\usepackage{booktabs}
\usepackage{tabularx}
\usepackage{array}
\usepackage{xcolor}
\usepackage[obeyspaces,spaces]{url}
\usepackage{xskak}
\usepackage{pgfplots}
\usepgfplotslibrary{groupplots}
\usetikzlibrary{positioning,arrows.meta,calc,fit}
\pgfplotsset{compat=1.18}
\usepackage{caption}
\usepackage{subcaption}
\usepackage{enumitem}
\usepackage[section]{placeins}
\usepackage[hidelinks]{hyperref}

% A generated value by its key: generated/values.tex sets them, \V{key} prints
% one, and a key nobody set stops the build.
\makeatletter
\newcommand{\setvalue}[2]{\expandafter\gdef\csname labvalue@#1\endcsname{#2}}
\newcommand{\V}[1]{\@ifundefined{labvalue@#1}%
  {\PackageError{lab}{No generated value named '#1'}{Run `lab paper` for this model again.}}%
  {\csname labvalue@#1\endcsname}}
\makeatother

% Chart colours. A run keeps its colour in every figure of every paper.
\definecolor{labBlue}{HTML}{2A78D6}
\definecolor{labOrange}{HTML}{EB6834}
\definecolor{labAqua}{HTML}{1BAF7A}
\definecolor{labYellow}{HTML}{EDA100}
\definecolor{labMagenta}{HTML}{E87BA4}
\definecolor{labGreen}{HTML}{008300}
\definecolor{labViolet}{HTML}{6250D6}
\definecolor{labGray}{HTML}{898781}
\definecolor{labInk}{HTML}{0B0B0B}
\definecolor{labSecondary}{HTML}{52514E}
\definecolor{labGrid}{HTML}{E1E0D9}
\definecolor{labAxis}{HTML}{C3C2B7}
\definecolor{labWash}{HTML}{CDE2FB}

\pgfplotsset{
  lab/.style={
    font=\sffamily\footnotesize,
    axis line style={labAxis},
    tick style={labAxis},
    tick label style={labSecondary},
    label style={labSecondary},
    title style={labInk, font=\sffamily\small\bfseries, align=left},
    grid style={labGrid, line width=0.3pt},
    major grid style={labGrid, line width=0.3pt},
    axis lines*=left,
    tick align=outside,
    clip=false,
    legend style={draw=none, fill=none, font=\sffamily\footnotesize, text=labSecondary,
      cells={anchor=west}, legend columns=-1, column sep=6pt, at={(0,1.03)}, anchor=south west},
    every axis plot/.append style={line width=1.1pt, line join=round, line cap=round},
    scaled ticks=false,
    /pgf/number format/1000 sep={,},
  },
  lab marks/.style={mark=*, mark size=1.9pt, mark options={draw=white, line width=0.6pt}},
}

\setchessboard{boardfontencoding=LSBC4, boardfontfamily=skaknew, blackfieldcolor=black!18,
  whitefieldcolor=white, setfontcolors, showmover=false, label=true, labelfontsize=6pt,
  boardfontsize=12pt}

% A move or line in UCI, a position in FEN (breakable at a slash), and a run's name.
\newcommand{\uci}[1]{\texttt{#1}}
\newcommand{\fen}[1]{{\footnotesize\path{#1}}}
\newcommand{\run}[1]{\texttt{#1}}
\newcolumntype{L}{>{\raggedright\arraybackslash}X}
\newcolumntype{R}{>{\raggedleft\arraybackslash}X}

\setlength{\parskip}{0.4em}
\setlength{\parindent}{0pt}
\captionsetup{font=small, labelfont=bf}

% from models/<model>/paper/, \input{../preamble} from papers/<name>/.
% Boards come from chessboard and xskak, plots from pgfplots. A paper prints no
% number by hand: `lab paper` writes every number, table and figure under the
% paper's generated/ folder, and paper.tex says what they mean.
\usepackage[LSBC4,T1]{fontenc}
\usepackage[utf8]{inputenc}
\usepackage{lmodern}
\usepackage[margin=1in]{geometry}
\usepackage{microtype}
\usepackage{amsmath}
\usepackage{booktabs}
\usepackage{tabularx}
\usepackage{array}
\usepackage{xcolor}
\usepackage[obeyspaces,spaces]{url}
\usepackage{xskak}
\usepackage{pgfplots}
\usepgfplotslibrary{groupplots}
\usetikzlibrary{positioning,arrows.meta,calc,fit}
\pgfplotsset{compat=1.18}
\usepackage{caption}
\usepackage{subcaption}
\usepackage{enumitem}
\usepackage[section]{placeins}
\usepackage[hidelinks]{hyperref}

% A generated value by its key: generated/values.tex sets them, \V{key} prints
% one, and a key nobody set stops the build.
\makeatletter
\newcommand{\setvalue}[2]{\expandafter\gdef\csname labvalue@#1\endcsname{#2}}
\newcommand{\V}[1]{\@ifundefined{labvalue@#1}%
  {\PackageError{lab}{No generated value named '#1'}{Run `lab paper` for this model again.}}%
  {\csname labvalue@#1\endcsname}}
\makeatother

% Chart colours. A run keeps its colour in every figure of every paper.
\definecolor{labBlue}{HTML}{2A78D6}
\definecolor{labOrange}{HTML}{EB6834}
\definecolor{labAqua}{HTML}{1BAF7A}
\definecolor{labYellow}{HTML}{EDA100}
\definecolor{labMagenta}{HTML}{E87BA4}
\definecolor{labGreen}{HTML}{008300}
\definecolor{labViolet}{HTML}{6250D6}
\definecolor{labGray}{HTML}{898781}
\definecolor{labInk}{HTML}{0B0B0B}
\definecolor{labSecondary}{HTML}{52514E}
\definecolor{labGrid}{HTML}{E1E0D9}
\definecolor{labAxis}{HTML}{C3C2B7}
\definecolor{labWash}{HTML}{CDE2FB}

\pgfplotsset{
  lab/.style={
    font=\sffamily\footnotesize,
    axis line style={labAxis},
    tick style={labAxis},
    tick label style={labSecondary},
    label style={labSecondary},
    title style={labInk, font=\sffamily\small\bfseries, align=left},
    grid style={labGrid, line width=0.3pt},
    major grid style={labGrid, line width=0.3pt},
    axis lines*=left,
    tick align=outside,
    clip=false,
    legend style={draw=none, fill=none, font=\sffamily\footnotesize, text=labSecondary,
      cells={anchor=west}, legend columns=-1, column sep=6pt, at={(0,1.03)}, anchor=south west},
    every axis plot/.append style={line width=1.1pt, line join=round, line cap=round},
    scaled ticks=false,
    /pgf/number format/1000 sep={,},
  },
  lab marks/.style={mark=*, mark size=1.9pt, mark options={draw=white, line width=0.6pt}},
}

\setchessboard{boardfontencoding=LSBC4, boardfontfamily=skaknew, blackfieldcolor=black!18,
  whitefieldcolor=white, setfontcolors, showmover=false, label=true, labelfontsize=6pt,
  boardfontsize=12pt}

% A move or line in UCI, a position in FEN (breakable at a slash), and a run's name.
\newcommand{\uci}[1]{\texttt{#1}}
\newcommand{\fen}[1]{{\footnotesize\path{#1}}}
\newcommand{\run}[1]{\texttt{#1}}
\newcolumntype{L}{>{\raggedright\arraybackslash}X}
\newcolumntype{R}{>{\raggedleft\arraybackslash}X}

\setlength{\parskip}{0.4em}
\setlength{\parindent}{0pt}
\captionsetup{font=small, labelfont=bf}

% from models/<model>/paper/, \input{../preamble} from papers/<name>/.
% Boards come from chessboard and xskak, plots from pgfplots. A paper prints no
% number by hand: `lab paper` writes every number, table and figure under the
% paper's generated/ folder, and paper.tex says what they mean.
\usepackage[LSBC4,T1]{fontenc}
\usepackage[utf8]{inputenc}
\usepackage{lmodern}
\usepackage[margin=1in]{geometry}
\usepackage{microtype}
\usepackage{amsmath}
\usepackage{booktabs}
\usepackage{tabularx}
\usepackage{array}
\usepackage{xcolor}
\usepackage[obeyspaces,spaces]{url}
\usepackage{xskak}
\usepackage{pgfplots}
\usepgfplotslibrary{groupplots}
\usetikzlibrary{positioning,arrows.meta,calc,fit}
\pgfplotsset{compat=1.18}
\usepackage{caption}
\usepackage{subcaption}
\usepackage{enumitem}
\usepackage[section]{placeins}
\usepackage[hidelinks]{hyperref}

% A generated value by its key: generated/values.tex sets them, \V{key} prints
% one, and a key nobody set stops the build.
\makeatletter
\newcommand{\setvalue}[2]{\expandafter\gdef\csname labvalue@#1\endcsname{#2}}
\newcommand{\V}[1]{\@ifundefined{labvalue@#1}%
  {\PackageError{lab}{No generated value named '#1'}{Run `lab paper` for this model again.}}%
  {\csname labvalue@#1\endcsname}}
\makeatother

% Chart colours. A run keeps its colour in every figure of every paper.
\definecolor{labBlue}{HTML}{2A78D6}
\definecolor{labOrange}{HTML}{EB6834}
\definecolor{labAqua}{HTML}{1BAF7A}
\definecolor{labYellow}{HTML}{EDA100}
\definecolor{labMagenta}{HTML}{E87BA4}
\definecolor{labGreen}{HTML}{008300}
\definecolor{labViolet}{HTML}{6250D6}
\definecolor{labGray}{HTML}{898781}
\definecolor{labInk}{HTML}{0B0B0B}
\definecolor{labSecondary}{HTML}{52514E}
\definecolor{labGrid}{HTML}{E1E0D9}
\definecolor{labAxis}{HTML}{C3C2B7}
\definecolor{labWash}{HTML}{CDE2FB}

\pgfplotsset{
  lab/.style={
    font=\sffamily\footnotesize,
    axis line style={labAxis},
    tick style={labAxis},
    tick label style={labSecondary},
    label style={labSecondary},
    title style={labInk, font=\sffamily\small\bfseries, align=left},
    grid style={labGrid, line width=0.3pt},
    major grid style={labGrid, line width=0.3pt},
    axis lines*=left,
    tick align=outside,
    clip=false,
    legend style={draw=none, fill=none, font=\sffamily\footnotesize, text=labSecondary,
      cells={anchor=west}, legend columns=-1, column sep=6pt, at={(0,1.03)}, anchor=south west},
    every axis plot/.append style={line width=1.1pt, line join=round, line cap=round},
    scaled ticks=false,
    /pgf/number format/1000 sep={,},
  },
  lab marks/.style={mark=*, mark size=1.9pt, mark options={draw=white, line width=0.6pt}},
}

\setchessboard{boardfontencoding=LSBC4, boardfontfamily=skaknew, blackfieldcolor=black!18,
  whitefieldcolor=white, setfontcolors, showmover=false, label=true, labelfontsize=6pt,
  boardfontsize=12pt}

% A move or line in UCI, a position in FEN (breakable at a slash), and a run's name.
\newcommand{\uci}[1]{\texttt{#1}}
\newcommand{\fen}[1]{{\footnotesize\path{#1}}}
\newcommand{\run}[1]{\texttt{#1}}
\newcolumntype{L}{>{\raggedright\arraybackslash}X}
\newcolumntype{R}{>{\raggedleft\arraybackslash}X}

\setlength{\parskip}{0.4em}
\setlength{\parindent}{0pt}
\captionsetup{font=small, labelfont=bf}

% Written by `lab paper think-time`; do not edit by hand.
\looktrue
\chancefalse
\sequencefalse
\contextgbdttrue
\pairstrue
\variantstrue
\shapestrue
\gamestrue
\featurestrue
\changestrue
\changehandtrue
\changeepochstrue
\changestarttrue
\changeattackstrue
\changegeometrytrue
\changeconvtrue
\changehybridtrue
\changelargerfalse
\moretrue
\shownfalse
\shownpairedfalse
\setvalue{band.1.floor}{0.7173}
\setvalue{band.1.name}{\textless{}1100}
\setvalue{band.5.floor}{0.8474}
\setvalue{band.5.name}{2000+}
\setvalue{bands.best-ahead.of}{all five bands}
\setvalue{bands.count}{five}
\setvalue{bands.lead.narrowest}{+0.0197}
\setvalue{bands.lead.narrowest.name}{1700-2000}
\setvalue{bands.lead.widest}{+0.0249}
\setvalue{bands.lead.widest.name}{1100-1400}
\setvalue{best.beside.closer}{53}
\setvalue{best.beside.correlation}{0.88}
\setvalue{best.close.median}{1.65}
\setvalue{best.close.p90}{3.57}
\setvalue{best.close.within.1.5}{42}
\setvalue{best.close.within.2}{64}
\setvalue{best.close.within.3}{85}
\setvalue{best.device}{NVIDIA A100-SXM4-40GB}
\setvalue{best.epochs.run}{12}
\setvalue{best.gain.of-lift}{15}
\setvalue{best.gain.percent}{3.4}
\setvalue{best.mae.factor}{1.84}
\setvalue{best.name}{\run{transformer-medium-hand-pretrained}}
\setvalue{best.parameters}{6{,}557{,}249}
\setvalue{best.r2}{0.3323}
\setvalue{best.reliability.groups}{20}
\setvalue{best.reliability.highest.happened}{24.0}
\setvalue{best.reliability.highest.said}{27.1}
\setvalue{best.reliability.lowest.happened}{1.6}
\setvalue{best.reliability.lowest.said}{1.4}
\setvalue{best.reliability.spread}{20}
\setvalue{best.reliability.widest-gap}{3.10}
\setvalue{best.rmse}{0.7772}
\setvalue{best.rows-a-second}{14{,}000}
\setvalue{best.score}{0.6089}
\setvalue{best.score.high}{0.6138}
\setvalue{best.score.low}{0.6042}
\setvalue{best.seconds}{9 min}
\setvalue{best.spearman}{0.5758}
\setvalue{challenger.name}{\run{transformer-medium-hand-pretrained}}
\setvalue{code.commit}{\texttt{36a0ad9}}
\setvalue{context.lift.share}{66}
\setvalue{cost.seconds}{106 min}
\setvalue{cost.seconds.networks}{103 min}
\setvalue{cost.seconds.others}{3 min}
\setvalue{data.fingerprint}{\texttt{340c68ed5c91}}
\setvalue{data.games.test}{2{,}944}
\setvalue{data.licence}{CC0 1.0}
\setvalue{data.month}{October 2026}
\setvalue{data.name}{\texttt{game-positions-2026-10}}
\setvalue{data.rows}{698{,}013}
\setvalue{data.rows.test}{68{,}579}
\setvalue{data.share.test}{10}
\setvalue{data.share.train}{85}
\setvalue{data.share.valid}{5}
\setvalue{describe.spent.mean}{12.1}
\setvalue{describe.spent.q25}{3}
\setvalue{describe.spent.q50}{7}
\setvalue{describe.spent.q95}{41}
\setvalue{describe.spent.q99}{81}
\setvalue{examples.count}{three}
\setvalue{floor.constant}{7.0}
\setvalue{floor.mae.factor}{2.16}
\setvalue{floor.rmse}{0.9511}
\setvalue{floor.score}{0.7724}
\setvalue{grid.rating-clock.-1100.-30.happened}{2.7}
\setvalue{grid.rating-clock.-1100.-30.rows}{47}
\setvalue{grid.rating-clock.-1100.300+.happened}{7.6}
\setvalue{grid.rating-clock.2000+.-30.happened}{1.8}
\setvalue{grid.rating-clock.2000+.-30.rows}{169}
\setvalue{grid.rating-clock.2000+.300+.happened}{6.8}
\setvalue{grid.rating-clock.gbm-hand.widest-off}{1.0}
\setvalue{grid.rating-clock.transformer-medium-hand-pretrained.widest-off}{1.0}
\setvalue{grid.smallest-cell}{30}
\setvalue{label.mean}{7.0}
\setvalue{ladder.ahead}{one}
\setvalue{ladder.at-last-epoch.of}{three of the four networks}
\setvalue{ladder.behind}{one}
\setvalue{ladder.count}{four}
\setvalue{ladder.level}{two}
\setvalue{look.date}{23 November 2026}
\setvalue{look.resamples}{200}
\setvalue{look.smallest-compared}{200}
\setvalue{look.test-rows}{68{,}579}
\setvalue{manifest.games-read}{1{,}003{,}668}
\setvalue{method.gbm.features.percent}{80}
\setvalue{method.gbm.leaves}{127}
\setvalue{method.gbm.patience}{100}
\setvalue{method.gbm.rate}{0.05}
\setvalue{method.gbm.rounds}{4{,}000}
\setvalue{method.gbm.rows.percent}{80}
\setvalue{method.resamples}{1{,}000}
\setvalue{more.gbm-hand.most}{four times}
\setvalue{networks.ahead.of}{nine of the twelve networks}
\setvalue{networks.patience}{three}
\setvalue{networks.warmup.percent}{5}
\setvalue{overall.gbm-hand.said}{6.8}
\setvalue{overall.happened}{7.0}
\setvalue{overall.transformer-medium-hand-pretrained.off}{$-$0.3}
\setvalue{overall.transformer-medium-hand-pretrained.said}{6.7}
\setvalue{pair.conv-medium-hand.transformer-medium-hand.gain}{+0.006249}
\setvalue{pair.conv-medium-hand.transformer-medium-hand.gain.high}{+0.007704}
\setvalue{pair.conv-medium-hand.transformer-medium-hand.gain.low}{+0.004890}
\setvalue{pair.conv-medium-hand.transformer-medium-hand.kept}{accepted}
\setvalue{pair.conv-medium-hand.transformer-medium-hand.valid.verdict}{better}
\setvalue{pair.conv-medium-hand.transformer-medium-hand.verdict}{significantly better}
\setvalue{pair.gbm-hand-double.gbm-hand.gain}{+0.002297}
\setvalue{pair.gbm-hand-double.gbm-hand.kept}{accepted}
\setvalue{pair.gbm-hand-double.gbm-hand.verdict}{significantly better}
\setvalue{pair.gbm-hand-half.gbm-hand.gain}{$-$0.001566}
\setvalue{pair.gbm-hand-half.gbm-hand.verdict}{significantly worse}
\setvalue{pair.gbm-hand-no-hanging.gbm-hand.gain}{$-$0.00965}
\setvalue{pair.gbm-hand-no-kings.gbm-hand.gain}{$-$0.002082}
\setvalue{pair.gbm-hand-no-mobility.gbm-hand.gain}{$-$0.008161}
\setvalue{pair.gbm-hand-no-pawns.gbm-hand.gain}{+0.0001036}
\setvalue{pair.gbm-hand-no-pawns.gbm-hand.verdict}{not significantly different}
\setvalue{pair.gbm-hand-quadruple.gbm-hand.gain}{+0.003630}
\setvalue{pair.gbm-hand-quadruple.gbm-hand.kept}{accepted}
\setvalue{pair.gbm-hand-quadruple.gbm-hand.verdict}{significantly better}
\setvalue{pair.gbm-hand-quarter.gbm-hand.gain}{$-$0.004568}
\setvalue{pair.gbm-hand-quarter.gbm-hand.verdict}{significantly worse}
\setvalue{pair.hybrid-medium-hand.transformer-medium-hand.gain}{$-$0.000796}
\setvalue{pair.hybrid-medium-hand.transformer-medium-hand.gain.high}{+0.000576}
\setvalue{pair.hybrid-medium-hand.transformer-medium-hand.gain.low}{$-$0.002055}
\setvalue{pair.hybrid-medium-hand.transformer-medium-hand.kept}{not accepted}
\setvalue{pair.hybrid-medium-hand.transformer-medium-hand.verdict}{not significantly different}
\setvalue{pair.transformer-medium-attacks.transformer-medium.gain}{+0.01787}
\setvalue{pair.transformer-medium-attacks.transformer-medium.gain.high}{+0.01934}
\setvalue{pair.transformer-medium-attacks.transformer-medium.gain.low}{+0.01650}
\setvalue{pair.transformer-medium-attacks.transformer-medium.kept}{accepted}
\setvalue{pair.transformer-medium-attacks.transformer-medium.valid.verdict}{better}
\setvalue{pair.transformer-medium-attacks.transformer-medium.verdict}{significantly better}
\setvalue{pair.transformer-medium-hand-attacks.transformer-medium-hand.gain}{+0.000576}
\setvalue{pair.transformer-medium-hand-attacks.transformer-medium-hand.gain.high}{+0.001023}
\setvalue{pair.transformer-medium-hand-attacks.transformer-medium-hand.gain.low}{+0.000133}
\setvalue{pair.transformer-medium-hand-attacks.transformer-medium-hand.kept}{not accepted}
\setvalue{pair.transformer-medium-hand-attacks.transformer-medium-hand.valid.verdict}{no better}
\setvalue{pair.transformer-medium-hand-attacks.transformer-medium-hand.verdict}{significantly better}
\setvalue{pair.transformer-medium-hand-double.transformer-medium-hand.gain}{+0.001551}
\setvalue{pair.transformer-medium-hand-double.transformer-medium-hand.kept}{accepted}
\setvalue{pair.transformer-medium-hand-double.transformer-medium-hand.verdict}{significantly better}
\setvalue{pair.transformer-medium-hand-half.transformer-medium-hand.gain}{$-$0.005738}
\setvalue{pair.transformer-medium-hand-long.transformer-medium-hand.gain}{+0.001414}
\setvalue{pair.transformer-medium-hand-long.transformer-medium-hand.gain.high}{+0.002486}
\setvalue{pair.transformer-medium-hand-long.transformer-medium-hand.gain.low}{+0.000382}
\setvalue{pair.transformer-medium-hand-long.transformer-medium-hand.kept}{accepted}
\setvalue{pair.transformer-medium-hand-long.transformer-medium-hand.valid.verdict}{better}
\setvalue{pair.transformer-medium-hand-long.transformer-medium-hand.verdict}{significantly better}
\setvalue{pair.transformer-medium-hand-long.transformer-medium-long.gain}{+0.00944}
\setvalue{pair.transformer-medium-hand-long.transformer-medium-long.gain.high}{+0.01114}
\setvalue{pair.transformer-medium-hand-long.transformer-medium-long.gain.low}{+0.00784}
\setvalue{pair.transformer-medium-hand-long.transformer-medium-long.verdict}{significantly better}
\setvalue{pair.transformer-medium-hand-no-geometry.transformer-medium-hand.gain}{$-$0.0002229}
\setvalue{pair.transformer-medium-hand-no-geometry.transformer-medium-hand.gain.high}{+0.0001984}
\setvalue{pair.transformer-medium-hand-no-geometry.transformer-medium-hand.gain.low}{$-$0.0006434}
\setvalue{pair.transformer-medium-hand-no-geometry.transformer-medium-hand.verdict}{not significantly different}
\setvalue{pair.transformer-medium-hand-pretrained.transformer-medium-hand.gain}{+0.01101}
\setvalue{pair.transformer-medium-hand-pretrained.transformer-medium-hand.gain.high}{+0.01270}
\setvalue{pair.transformer-medium-hand-pretrained.transformer-medium-hand.gain.low}{+0.00928}
\setvalue{pair.transformer-medium-hand-pretrained.transformer-medium-hand.kept}{accepted}
\setvalue{pair.transformer-medium-hand-pretrained.transformer-medium-hand.valid.verdict}{better}
\setvalue{pair.transformer-medium-hand-pretrained.transformer-medium-hand.verdict}{significantly better}
\setvalue{pair.transformer-medium-hand-quadruple.transformer-medium-hand.gain}{+0.002044}
\setvalue{pair.transformer-medium-hand-quadruple.transformer-medium-hand.kept}{accepted}
\setvalue{pair.transformer-medium-hand-quadruple.transformer-medium-hand.verdict}{significantly better}
\setvalue{pair.transformer-medium-hand-quarter.transformer-medium-hand.gain}{$-$0.01256}
\setvalue{pair.transformer-medium-hand.transformer-medium.gain}{+0.01276}
\setvalue{pair.transformer-medium-hand.transformer-medium.gain.high}{+0.01434}
\setvalue{pair.transformer-medium-hand.transformer-medium.gain.low}{+0.01124}
\setvalue{pair.transformer-medium-hand.transformer-medium.kept}{accepted}
\setvalue{pair.transformer-medium-hand.transformer-medium.valid.verdict}{better}
\setvalue{pair.transformer-medium-hand.transformer-medium.verdict}{significantly better}
\setvalue{pair.transformer-medium-long.transformer-medium.gain}{+0.004728}
\setvalue{pair.transformer-medium-long.transformer-medium.gain.high}{+0.005666}
\setvalue{pair.transformer-medium-long.transformer-medium.gain.low}{+0.003745}
\setvalue{pair.transformer-medium-long.transformer-medium.kept}{accepted}
\setvalue{pair.transformer-medium-long.transformer-medium.valid.verdict}{better}
\setvalue{pair.transformer-medium-long.transformer-medium.verdict}{significantly better}
\setvalue{pair.transformer-medium-pretrained.transformer-medium.gain}{+0.02286}
\setvalue{pair.transformer-medium-pretrained.transformer-medium.gain.high}{+0.02493}
\setvalue{pair.transformer-medium-pretrained.transformer-medium.gain.low}{+0.02099}
\setvalue{pair.transformer-medium-pretrained.transformer-medium.kept}{accepted}
\setvalue{pair.transformer-medium-pretrained.transformer-medium.valid.verdict}{better}
\setvalue{pair.transformer-medium-pretrained.transformer-medium.verdict}{significantly better}
\setvalue{recipe.games}{30{,}000}
\setvalue{recipe.min-ply}{11}
\setvalue{recipe.per-game}{24}
\setvalue{reference.close.within.1.5}{40}
\setvalue{reference.close.within.2}{62}
\setvalue{reference.close.within.3}{84}
\setvalue{reference.features}{24}
\setvalue{reference.leans.1.feature}{mover's clock}
\setvalue{reference.leans.1.share}{21}
\setvalue{reference.leans.2.feature}{legal moves}
\setvalue{reference.leans.2.share}{11}
\setvalue{reference.leans.3.feature}{increment}
\setvalue{reference.leans.3.share}{8}
\setvalue{reference.mae.factor}{1.88}
\setvalue{reference.r2}{0.2939}
\setvalue{reference.reliability.highest.said}{23.6}
\setvalue{reference.reliability.lowest.said}{1.6}
\setvalue{reference.reliability.widest-gap}{1.20}
\setvalue{reference.rmse}{0.7992}
\setvalue{reference.score}{0.6306}
\setvalue{reference.spearman}{0.5322}
\setvalue{reference.trees}{709}
\setvalue{results.date}{23 November 2026}
\setvalue{run.conv-medium-hand.lr}{0.002}
\setvalue{run.conv-medium-hand.parameters.millions}{6.9}
\setvalue{run.conv-medium-hand.score}{0.6137}
\setvalue{run.conv-small.best-epoch}{8}
\setvalue{run.conv-small.epochs}{8}
\setvalue{run.conv-small.gain}{+0.001397}
\setvalue{run.conv-small.gain.high}{+0.003967}
\setvalue{run.conv-small.gain.low}{$-$0.001706}
\setvalue{run.conv-small.parameters.millions}{1.0}
\setvalue{run.gbm-context.gain}{$-$0.04885}
\setvalue{run.gbm-context.r2}{0.1945}
\setvalue{run.gbm-context.score}{0.6794}
\setvalue{run.gbm-context.spearman}{0.4265}
\setvalue{run.gbm-hand-quadruple.score}{0.6270}
\setvalue{run.gbm-hand-quadruple.trees}{744}
\setvalue{run.gbm-hand-quarter.score}{0.6352}
\setvalue{run.hybrid-medium-hand.convs}{2}
\setvalue{run.hybrid-medium-hand.depth}{5}
\setvalue{run.hybrid-medium-hand.parameters.millions}{6.6}
\setvalue{run.transformer-medium-attacks.gain.high}{+0.01746}
\setvalue{run.transformer-medium-attacks.gain.low}{+0.01364}
\setvalue{run.transformer-medium-attacks.score}{0.6148}
\setvalue{run.transformer-medium-hand-attacks.gain.high}{+0.01270}
\setvalue{run.transformer-medium-hand-attacks.gain.low}{+0.00963}
\setvalue{run.transformer-medium-hand-attacks.score}{0.6194}
\setvalue{run.transformer-medium-hand-double.best-epoch}{6}
\setvalue{run.transformer-medium-hand-double.epochs}{6}
\setvalue{run.transformer-medium-hand-long.best-epoch}{21}
\setvalue{run.transformer-medium-hand-pretrained.best-epoch}{10}
\setvalue{run.transformer-medium-hand-pretrained.epochs}{12}
\setvalue{run.transformer-medium-hand-pretrained.over-twin}{+0.000900}
\setvalue{run.transformer-medium-hand-pretrained.score}{0.6089}
\setvalue{run.transformer-medium-hand-quadruple.best-epoch}{3}
\setvalue{run.transformer-medium-hand-quadruple.epochs}{3}
\setvalue{run.transformer-medium-hand-quadruple.score}{0.6179}
\setvalue{run.transformer-medium-hand-quarter.score}{0.6325}
\setvalue{run.transformer-medium-hand.best-epoch}{11}
\setvalue{run.transformer-medium-hand.epochs}{12}
\setvalue{run.transformer-medium-hand.gain}{+0.01063}
\setvalue{run.transformer-medium-hand.gain.high}{+0.01214}
\setvalue{run.transformer-medium-hand.gain.low}{+0.00907}
\setvalue{run.transformer-medium-hand.lr}{0.0006}
\setvalue{run.transformer-medium-hand.parameters.millions}{6.6}
\setvalue{run.transformer-medium-hand.score}{0.6200}
\setvalue{run.transformer-medium-long.best-epoch}{22}
\setvalue{run.transformer-medium-long.epochs}{24}
\setvalue{run.transformer-medium-long.gain}{+0.002601}
\setvalue{run.transformer-medium-long.gain.high}{+0.004526}
\setvalue{run.transformer-medium-long.gain.low}{+0.000209}
\setvalue{run.transformer-medium-pretrained.best-epoch}{10}
\setvalue{run.transformer-medium-pretrained.score}{0.6098}
\setvalue{run.transformer-medium.best-epoch}{12}
\setvalue{run.transformer-medium.epochs}{12}
\setvalue{run.transformer-medium.gain}{$-$0.002127}
\setvalue{run.transformer-medium.gain.high}{$-$0.000343}
\setvalue{run.transformer-medium.gain.low}{$-$0.004274}
\setvalue{run.transformer-medium.parameters.millions}{6.6}
\setvalue{run.transformer-small.best-epoch}{8}
\setvalue{run.transformer-small.epochs}{8}
\setvalue{run.transformer-small.gain}{+0.001931}
\setvalue{run.transformer-small.gain.high}{+0.003677}
\setvalue{run.transformer-small.gain.low}{$-$0.000207}
\setvalue{run.transformer-small.parameters.millions}{0.9}
\setvalue{runs.count}{17}
\setvalue{slice.band.-1100.happened}{7.5}
\setvalue{slice.band.2000+.happened}{6.5}
\setvalue{slice.band.ahead.of}{all five groups}
\setvalue{slice.captures.0.gain}{+0.00708}
\setvalue{slice.captures.0.happened}{5.2}
\setvalue{slice.captures.4+.gain}{+0.03248}
\setvalue{slice.captures.4+.happened}{8.8}
\setvalue{slice.check.no.happened}{7.3}
\setvalue{slice.check.yes.happened}{4.4}
\setvalue{slice.clock.-30.happened}{2.7}
\setvalue{slice.clock.30-60.happened}{4.6}
\setvalue{slice.clock.30-60.rows}{777}
\setvalue{slice.clock.300-480.happened}{8.1}
\setvalue{slice.clock.480+.happened}{6.6}
\setvalue{slice.clock.60-120.happened}{6.2}
\setvalue{slice.clock.60-120.rows}{1{,}456}
\setvalue{slice.clock.ahead.groups}{\textless{}30, 120-300, 300-480 and 480+}
\setvalue{slice.clock.ahead.of}{four of the six groups}
\setvalue{slice.clock.challenger.widest-off}{0.5}
\setvalue{slice.clock.reference.widest-off}{0.3}
\setvalue{slice.hanging.0.happened}{6.8}
\setvalue{slice.hanging.3+.happened}{7.8}
\setvalue{slice.moves.-5.happened}{4.1}
\setvalue{slice.moves.45+.happened}{9.6}
\setvalue{slice.moves.5-15.happened}{3.9}
\setvalue{slice.moves.ahead.groups}{\textless{}5, 15-25, 25-35, 35-45 and 45+}
\setvalue{slice.moves.gain.least}{+0.00397}
\setvalue{slice.moves.gain.least.group}{5-15}
\setvalue{slice.pieces.-3.gain}{+0.00836}
\setvalue{slice.pieces.-3.happened}{3.0}
\setvalue{slice.pieces.9-12.happened}{8.3}
\setvalue{slice.pieces.ahead.of}{all five groups}
\setvalue{slice.ply.-20.happened}{5.7}
\setvalue{slice.ply.30-40.happened}{8.7}
\setvalue{slice.ply.80+.gain}{+0.01326}
\setvalue{slice.ply.80+.happened}{3.4}
\setvalue{slice.ply.ahead.of}{all six groups}
\setvalue{slice.seconds.-2.challenger.said}{3.0}
\setvalue{slice.seconds.-2.reference.said}{3.6}
\setvalue{slice.seconds.5-10.gain}{$-$0.05022}
\setvalue{slice.seconds.80+.challenger.said}{17.4}
\setvalue{slice.seconds.80+.happened}{110.4}
\setvalue{slice.seconds.80+.reference.said}{15.2}
\setvalue{slice.seconds.ahead.groups}{\textless{}2, 2-5, 20-40, 40-80 and 80+}
\setvalue{slice.seconds.ahead.of}{five of the seven groups}
\setvalue{slice.seconds.behind.groups}{5-10}
\setvalue{slice.seconds.behind.of}{one of the seven groups}
\setvalue{slice.win.-10.challenger.off}{$-$0.5}
\setvalue{slice.win.-10.happened}{6.1}
\setvalue{slice.win.-10.reference.off}{$-$0.7}
\setvalue{slice.win.10-20.happened}{7.8}
\setvalue{slice.win.20-40.happened}{7.9}
\setvalue{slice.win.40-60.challenger.off}{$-$0.6}
\setvalue{slice.win.40-60.happened}{6.7}
\setvalue{slice.win.60-80.happened}{8.2}
\setvalue{slice.win.80-90.happened}{8.1}
\setvalue{slice.win.90+.challenger.off}{+0.3}
\setvalue{slice.win.90+.happened}{6.4}
\setvalue{slice.win.90+.reference.off}{+0.5}
\setvalue{slice.win.ahead.of}{all seven groups}
\setvalue{study.features.gbm-hand.most.gain}{$-$0.01410}
\setvalue{study.features.gbm-hand.most.group}{material}
\setvalue{study.features.gbm-hand.worse.of}{four of the five ablations}
\setvalue{study.games.transformer-medium-hand.worse.of}{both ablations}
\setvalue{task.context.count}{eleven}
\setvalue{task.hand.count}{13}
\setvalue{task.marks.count}{three}
\setvalue{task.reference}{\run{gbm-hand}}


\begin{document}

\papertitle{Predicting Move Times in Human Chess Games}
  {Results as of \V{results.date}\quad\textperiodcentered\quad dataset \V{data.name},
   build \V{data.fingerprint}\quad\textperiodcentered\quad code \V{code.commit}}

\begin{abstract}
How long a chess player thinks about a move depends on the position, on the
clock and on the player. We train models that predict the time spent on a
move in a rated rapid game from the position, the opponent's preceding move,
both players' ratings, both clocks and the time control, without an engine
evaluation. The target is the logarithm of one plus the number of seconds,
and the primary metric is the mean absolute error (MAE) of this logarithm, so
that an error corresponds to a factor. The dataset contains \V{data.rows}
positions from \V{recipe.games} rated rapid games played on Lichess, split by
game. We compare a constant baseline, gradient-boosted decision trees (GBDT)
on hand-crafted features, and convolutional and transformer networks that
read the board as 64 square tokens. The best of \V{runs.count} models is a
transformer that also receives the hand-crafted features and whose encoder is
pretrained on engine evaluations. On \V{data.rows.test} test positions its
MAE is \V{best.score} (95\% CI \V{best.score.low} to \V{best.score.high}), an
error factor of \V{best.mae.factor}, compared with \V{reference.score}
(\V{reference.mae.factor}) for the GBDT baseline and \V{floor.score}
(\V{floor.mae.factor}) for the constant baseline. GBDT on the scalar context
features alone, which are the ratings, the clocks, the time control, the ply
and the castling rights, achieves \V{context.lift.share}\% of the baseline's
improvement over the constant. Pretraining, an attack-aware attention bias, a
residual convolutional encoder, a longer training schedule and more training
data each improve a network significantly. The coefficient of determination
of the best model is \V{best.r2}. Its predictions are accurate on average
over their whole range, but they do not anticipate individual long thinks: a
single position does not determine most of the variation in the time a move
takes.
\end{abstract}

\section{Introduction}
\label{sec:intro}

A chess engine evaluates a position without regard to the clock, and the
clock shows the time that remains without regard to the position. The time a
player spends on a move depends on both. Eight seconds are a short time in a
position with three captures to consider, and a long time for a forced
recapture.

This work supports a chess coaching application that gives feedback on how a
player uses the clock in rapid games: a move that was played too quickly, or
one that took longer than the position required. Such feedback needs a
reference. A fixed number of seconds cannot serve as one, because it ignores
both the position and the clock. The time that players of the same rating
usually take in the same position with the same time remaining can, and a
model of this time provides the reference.

We study the following prediction task. Before a move is played, given the
position, the opponent's preceding move, both players' ratings, both clock
times and the time control, predict the time that the player will spend on
the move. No model receives an engine evaluation. A player does not know the
evaluation, and the reference is meant to describe what players do with the
information that is available to them.

We evaluate a fixed sequence of models on \V{data.rows} positions from rated
rapid games on Lichess: a constant baseline, GBDT on scalar context features,
GBDT with hand-crafted features of the position, and neural networks that
read the board directly, with and without the hand-crafted features. Because
the sequence is fixed, each step measures the contribution of one component.
Our main findings are the following.
\begin{itemize}
\item \textbf{The context accounts for most of what is predicted.} GBDT on
the scalar context features alone (the ratings, the clocks, the time control,
the ply and the castling rights), without the board, has an MAE of
\V{run.gbm-context.score}. It achieves \V{context.lift.share}\% of the
improvement of the GBDT baseline (\V{reference.score}) over the constant
baseline (\V{floor.score}).
\item \textbf{A network that reads the board improves on GBDT by a small
margin.} With the hand-crafted features, the medium transformer is
significantly better than the GBDT baseline (gain
\V{run.transformer-medium-hand.gain}; 95\% CI
\V{run.transformer-medium-hand.gain.low} to
\V{run.transformer-medium-hand.gain.high}). The best model's gain is
\V{best.gain.percent}\% of the baseline's MAE (Section~\ref{sec:results}).
\ifchanges
\item \textbf{Several modifications of the network are
accepted.}\ifchangestart{} Initialising the transformer's encoder from a
model trained on engine evaluations gives the largest paired gain,
\V{pair.transformer-medium-pretrained.transformer-medium.gain} without the
hand-crafted features.\fi\ifchangeattacks{} An attention bias computed from
the attacks between squares gives
\V{pair.transformer-medium-attacks.transformer-medium.gain}, more than the
hand-crafted features do
(\V{pair.transformer-medium-hand.transformer-medium.gain}).\fi\ifchangeconv{}
A residual convolutional encoder is significantly better than a transformer
encoder of similar size.\fi\ifchangeepochs{} A schedule of twice the length
is accepted for both medium transformers.\fi{}
(Section~\ref{sec:ablations}.)
\fi
\ifmore
\item \textbf{More training data helps both model classes.} With four times
the training games, the gain is \V{pair.gbm-hand-quadruple.gbm-hand.gain} for
GBDT and
\V{pair.transformer-medium-hand-quadruple.transformer-medium-hand.gain} for
the transformer (Section~\ref{sec:training-size}).
\fi
\item \textbf{Most of the variation is not predictable from one position.}
The coefficient of determination (R\textsuperscript{2}) of the best model
is \V{best.r2}\iflook, and \V{best.close.within.2}\% of its predictions are
within a factor of two of the time spent\fi.
\iflook
\item \textbf{The predictions are accurate on average and do not anticipate
long thinks.} When the test positions are divided into
\V{best.reliability.groups} bins by the predicted time, the mean prediction
and the typical time spent, both computed on the logarithmic scale
(Section~\ref{sec:protocol}), differ by at most
\V{best.reliability.widest-gap} seconds in any bin. The highest bin has a
mean prediction of
\V{best.reliability.highest.said} seconds, whereas one move in twenty takes
more than \V{describe.spent.q95} seconds (Section~\ref{sec:error}).
\fi
\end{itemize}

\section{Related Work}
\label{sec:related}

\paragraph{Models of human play.}
Maia~\cite{maia} and Maia-2~\cite{maia2} predict the move that a player of a
given rating plays in a position. They model the choice of the move and not
the time it takes. Allie~\cite{allie} is a transformer trained on sequences
of moves from human games that predicts, in addition to the move, how long
the player ponders and whether the player resigns. It reads the game as a
sequence. We study a narrower setting: one position with its clocks, without
the preceding moves and without an engine. We measure how much of the time
spent can be predicted in this setting, and how much each input and each
model component contributes relative to a feature-based baseline on the same
test games.

\paragraph{Time and errors.}
Anderson et al.~\cite{humanerror} relate the errors of human players to the
difficulty of the position, the player's skill and the time available, and
find that the difficulty of the position predicts errors best. A model of
blunder risk on the positions used here is described in a companion paper of
this repository. Its error analysis shows that the blunder rate increases
with the time spent on the move, which is one motivation for modelling that
time directly.

\paragraph{Pretraining on engine evaluations.}
Ruoss et al.~\cite{searchless} show that a large transformer trained by
supervised learning on engine evaluations plays strongly without search. We
use a far smaller model of this kind only as a source of pretrained weights
(Section~\ref{sec:ablations}).

\section{Task and Data}
\label{sec:data}

\subsection{Task}

An example is one position of a game, immediately before a move is played.
The input is the position, the opponent's preceding move, both players'
ratings, both clock times, the time control and the ply. The target is the
time that the player spent on the move, on a logarithmic scale.
Table~\ref{tab:brief} summarises the task.

\begin{table}[htbp]\centering\small
\begin{tabularx}{\textwidth}{lL}
\toprule
Input & The position before the move, the opponent's preceding move, both
players' ratings, both clock times, the time control and the ply. No engine
evaluation. \\
Target & The time spent on the move, as $\log(1 + \text{seconds})$. \\
Data & Rated rapid games from the Lichess database (\V{data.licence}) for
\V{data.month} that carry clock times and engine evaluations:
\V{recipe.games} games and up to \V{recipe.per-game} positions per game. \\
Split & By game: \V{data.share.train}\% training, \V{data.share.valid}\%
validation, \V{data.share.test}\% test. \\
Primary metric & MAE of the target on the test set (lower is better). An
error of 0.69 corresponds to a factor of two. \\
\bottomrule
\end{tabularx}
\caption{Summary of the task.}
\label{tab:brief}
\end{table}

\paragraph{Definition of the time spent.}
Lichess records the time remaining on a player's clock after each of the
player's moves, in whole seconds. Let $t_{\mathrm{before}}$ be the time on
the mover's clock after the mover's previous move, $t_{\mathrm{after}}$ the
time on it after the move in question, and $i$ the increment that the time
control adds after every move. The time spent on the move is
\[
s = \max(0,\; t_{\mathrm{before}} - t_{\mathrm{after}} + i),
\]
and the target is $y = \log(1 + s)$, with the natural logarithm. Because the
clocks are recorded in whole seconds, $s$ is a whole number, and a move that
is played immediately has $s = 0$ or $s = 1$.

\subsection{Dataset}

The source is the Lichess database of rated games~\cite{lichessdb} for
\V{data.month}. A game is included when it is rated, both players have a
rating and neither is a bot, the time control is rapid by Lichess's
definition (an estimated duration of 8 to 25 minutes, where the estimate is
the initial time plus 40 times the increment), the game ended on the board
or on time, and its record carries engine evaluations and clock times. The
dataset is shared with the blunder-risk task, whose label is computed from
the evaluations. This is why games without evaluations are excluded although
no model in this paper receives an evaluation. Lichess computes evaluations
for a game only when a user requests a computer analysis of it, so the
analysed games are a self-selected subset of all games.

The monthly file is read from its beginning until \V{recipe.games} games have
contributed positions, and \V{manifest.games-read} games were read to find
them. The dataset therefore consists of the first qualifying games of the
month in file order. It is not a random sample of the month.

From each game we take up to \V{recipe.per-game} positions, selected among
the eligible plies by a hash of the game identifier and the ply. A ply is
eligible from ply \V{recipe.min-ply} on, that is after each side's first five
moves, the opening cut used by Maia~\cite{maia}, provided that the
evaluations before and after the move and the clock times around the move are
all available. The resulting dataset has \V{data.rows} positions
(Table~\ref{tab:data}).

The split is determined by a hash of the game identifier, so all positions
of a game fall in the same set, and the test set contains positions from
\V{data.games.test} games. The same grouping is used for the bootstrap
(Section~\ref{sec:protocol}). The dataset is not stored in the repository.
Its recipe is, and a build of the dataset is identified by a fingerprint, a
hash of its contents. All models in this paper were trained and evaluated on
build \V{data.fingerprint}\ifmore; the models of
Section~\ref{sec:training-size} that use more training games were trained on
a larger build and evaluated on the validation and test sets of this
one\fi. The evaluation code refuses to rank models evaluated on different
builds.

\begin{table}[htbp]\centering\small
\block{data}
\caption{Dataset statistics by split. The mean target is converted to
seconds: it is the typical time defined in Section~\ref{sec:protocol}.}
\label{tab:data}
\end{table}

\iflook
\paragraph{Distribution of the target.}
The time spent has a long right tail (Figure~\ref{fig:label},
Table~\ref{tab:describe}). Half of all moves take \V{describe.spent.q50}
seconds or less and a quarter take \V{describe.spent.q25} seconds or less,
while one move in twenty takes more than \V{describe.spent.q95} seconds and
one in a hundred more than \V{describe.spent.q99}. This is the reason for the
logarithmic target. On a linear scale, the error of a model would be
dominated by a small number of very long thinks. On the logarithmic scale,
the difference between two and four seconds counts as much as the difference
between twenty and forty.

\begin{figure}[htbp]\centering
\block{label}
\caption{Distribution of the time spent on a move over all positions of the
dataset, on the scale of the target.}
\label{fig:label}
\end{figure}

\begin{table}[htbp]\centering\small
\block{describe}
\caption{Summary statistics of selected variables over the whole dataset:
the mean, five quantiles, the minimum and the maximum. Clock times and the
time spent are in seconds. The time spent is the target and is not an
input.}
\label{tab:describe}
\end{table}
\fi

\section{Methods}
\label{sec:methods}

\subsection{Input representation}
\label{sec:input}

% Input representation: identical for every task, so each model paper inputs
% this file under its own heading and then adds what is specific to its task.
% The prose follows docs/writing.md.
All tasks in this project share one input representation, so two tasks differ
in their target and in the information a model is given, and not in how a
position is encoded.

\paragraph{Square tokens.}
A position is encoded as 64 tokens, one per square. Token 0 denotes an empty
square, tokens 1 to 6 the pawn, knight, bishop, rook, queen and king of the
side to move, which we call the \emph{mover}, and tokens 7 to 12 the
opponent's pieces. When Black is to move, the board is flipped vertically and
the colours are swapped, so that the mover always plays up the board. A
position with White to move and the corresponding colour-reversed position
with Black to move therefore have the same encoding, and a model does not
have to learn the two cases separately.

\paragraph{Auxiliary input planes.}
Information about individual squares is supplied as binary planes over the 64
squares, which a network receives together with the tokens. This task uses
\V{task.marks.count} planes (Table~\ref{tab:marks}). The first marks the en
passant target square, and the others mark the origin and destination squares
of moves. A plane identifies the squares of a move and carries no other
information about it, such as whether the move is a capture or a check.

\begin{table}[htbp]\centering\small
\block{marks}
\caption{Auxiliary input planes. Squares are given from the mover's point of
view, after the board has been flipped.}
\label{tab:marks}
\end{table}

\ifsequence
\paragraph{Sequences of positions.}
An example of this task is not a single position. It is the sequence of one
player's moves in a game, and each element of the sequence is a position
encoded as above: the board before the move, with the opponent's preceding
move and the move played marked in the auxiliary planes. Scalar inputs
therefore come at two levels, one for each move and one for the game.

\paragraph{Per-move features.}
Each position of the sequence has \V{task.steps.count} scalar features of its
own (Table~\ref{tab:steps}): the four castling rights and the task's
information about that move. A network receives them together with the
position they belong to.

\begin{table}[htbp]\centering\small
\block{steps}
\caption{Per-move features: the scalar inputs attached to each position of a
sequence.}
\label{tab:steps}
\end{table}

\paragraph{Scalar context features.}
Information about the game as a whole is given as \V{task.context.count}
scalar context features (Table~\ref{tab:reads}). Every model receives them,
including the gradient-boosted trees, which never see a board. For a network,
both kinds of scalar feature are standardised with the mean and standard
deviation of the training set.

\begin{table}[htbp]\centering\small
\block{reads}
\caption{Scalar context features: the game-level inputs given to every
model.}
\label{tab:reads}
\end{table}

\paragraph{Hand-crafted features.}
Gradient-boosted trees cannot take a sequence of boards as input, so we also
compute \V{task.hand.count} aggregate features over the player's moves
(Table~\ref{tab:hand} in the appendix): how much win probability the moves
lost on average, at worst and how often beyond fixed thresholds, how the clock
was used, and how many moves were available. These features have two roles.
They are the inputs of the GBDT baseline, and so define the level that a
network reading every position has to exceed. They can also be given to a
network as additional context features, which tests whether the network
already extracts the same information from the sequence.

\else
\paragraph{Scalar context features.}
Information that does not belong to a square is given as scalar context
features: the four castling rights and the scalar inputs specific to the
task. This task uses \V{task.context.count} (Table~\ref{tab:reads}). Every
model receives them, including the gradient-boosted trees, which never see
the board. For a network they are standardised with the mean and standard
deviation of the training set.

\begin{table}[htbp]\centering\small
\block{reads}
\caption{Scalar context features: the inputs given to every model in addition
to the board.}
\label{tab:reads}
\end{table}

\paragraph{Hand-crafted features.}
Gradient-boosted trees cannot take 64 tokens as input, so we also compute
\V{task.hand.count} hand-crafted features of the position
(Table~\ref{tab:hand} in the appendix): material, mobility, available checks
and captures, hanging pieces, passed pawns and attacks near each king, and,
where the task gives a line of moves, features of that line. These features
have two roles. They are the inputs of the GBDT baseline, and so define the
level that a network reading the board has to exceed. They can also be given
to a network as additional context features, which tests whether the network
already extracts the same information from the board.

\fi

\iflook
\begin{figure}[htbp]\centering
\block{encoding}
\ifsequence
\caption{One position of one example, shown twice: the position before the
player's most costly move of the game. Left: as played, with arrows for the
opponent's preceding move (grey) and the player's move (blue). Right: as
encoded, oriented so that the mover plays up the board, with the squares set
in an auxiliary plane shaded grey for the opponent's move and blue for the
player's move. Every position of the sequence is encoded in this way.}
\else
\caption{One example from the dataset, shown twice. Left: as played, with
arrows for the opponent's preceding move (grey) and the move the example
concerns (blue). Right: as encoded, oriented so that the mover plays up the
board, with the squares set in an auxiliary plane shaded grey for the
opponent's move and, where the task marks a line of moves, blue for the
mover's moves and orange for the replies.}
\fi
\label{fig:encoding}
\end{figure}
\fi


\paragraph{Inputs of this task.}
The auxiliary planes mark the en passant square and the origin and
destination squares of the opponent's preceding move. Besides the castling
rights, the scalar context features are the two ratings, the increment, the
ply, and three times, each as the logarithm of one plus the number of
seconds: the time remaining on the mover's clock, the time remaining on the
opponent's clock and the initial time of the time control. No model receives
an engine evaluation, although the dataset contains one for every position.
No input describes the move that is played. The move, the time spent on it
and the evaluations appear in the error analysis
(Section~\ref{sec:error}).\iflook{} In Figure~\ref{fig:encoding} the move
played is drawn on the left board for the reader, and it is not marked on the
encoded board.\fi

\subsection{Models and training}
\label{sec:models}

% Models and training: the same sequence of models for every task, so each
% model paper inputs this file under its own heading.
% The prose follows docs/writing.md.
Every task is given the same sequence of models, so that each step of the
sequence measures the contribution of one component. A model is fitted on the
training set, selected on the validation set and scored once on the test set.
Table~\ref{tab:runs} lists the models of this task.

\paragraph{Constant baseline.}
\run{constant} predicts the same value for every example: the mean of the
training targets (for a binary target, the base rate), or their median when
the primary metric is an absolute error. A useful model has to improve on it.

\paragraph{Gradient-boosted decision trees.}
\ifcontextgbdt
\run{gbm-context} is a gradient-boosted decision tree model (GBDT), trained
with LightGBM~\cite{lightgbm} on the scalar context features alone. It never
sees the board, and so measures how much of the target the context already
determines. \run{gbm-hand} adds the hand-crafted features.
\else
\run{gbm-hand} is a gradient-boosted decision tree model (GBDT), trained with
LightGBM~\cite{lightgbm} on the scalar context features and the hand-crafted
features.
\fi
We refer to \run{gbm-hand} as \emph{the GBDT baseline}: every other model is
compared with it on the same test examples. GBDT is trained for at most
\V{method.gbm.rounds} boosting rounds with up to \V{method.gbm.leaves} leaves
per tree and a learning rate of \V{method.gbm.rate}. Each tree is fitted on a
random \V{method.gbm.rows.percent}\% of the examples and
\V{method.gbm.features.percent}\% of the features, and boosting stops when the
validation score has not improved for \V{method.gbm.patience} rounds.

\paragraph{Neural networks.}
A network consists of an encoder, which maps the square tokens, the auxiliary
planes and the context features to a vector, and a prediction head, a
two-layer perceptron that maps this vector to the output. We use two encoder
architectures, so that the results show how much the architecture matters for
the task.

The \emph{transformer} encoder~\cite{attention} treats each square as a token.
The embedding of a token is the sum of a learned embedding of its piece, a
learned embedding of its square and a linear projection of the square's
auxiliary planes. The context features are projected by a small perceptron
and enter as one additional token. Every token attends to every other token
in every layer, so distant squares interact from the first layer on. To give
attention the geometry of the board, each attention head adds a learned bias
that depends on the file and rank offset between two
squares~\cite{relative}: files, ranks, diagonals and knight moves all
correspond to fixed offsets. The layers use pre-layer normalisation, and the
encoder outputs the context token concatenated with the mean of the 64 square
tokens.

The \emph{residual convolutional network} (ResNet) is a stack of residual
blocks of $3\times3$ convolutions over the $8\times8$ board, the architecture
that AlphaZero introduced for chess~\cite{alphazero}. Its input is one binary
plane per token value together with the auxiliary planes. A linear projection
of the context features is added at every square after the first convolution,
and the encoder outputs the mean and the maximum over the squares of the
final feature map. Unlike attention, a convolution relates only neighbouring
squares, so the receptive field grows with depth.

\ifsequence
\paragraph{Sequence model.}
An example of this task is a sequence of positions, so a network has two
stages below its prediction head. One of the two encoders above is applied to
every position of the sequence, with the position's per-move features in
place of the context features, and produces one vector per position. A
two-layer transformer then operates on these vectors in the order in which
the moves were played: each vector is a token, a learned positional embedding
is added to it, and the game-level context features enter as one additional
token. The prediction head receives that token concatenated with the mean of
the position tokens. Sequences in a batch are padded to the length of the
longest, and padded elements are masked: they are not encoded and no token
attends to them, so a prediction does not depend on the other examples in
the batch.

Model sizes are hyperparameters. \run{conv-small} is small enough to train
on a CPU. An example consists of tens of positions, so
\run{transformer-small} and \run{transformer-medium} are trained on a GPU.
\run{transformer-medium-hand} is \run{transformer-medium} with the
hand-crafted features added to its context features.
Table~\ref{tab:settings} in the appendix lists the hyperparameters.
\else
Model sizes are hyperparameters. \run{conv-small} and \run{transformer-small}
are small enough to train on a CPU. \run{transformer-medium} is trained on a
GPU, and \run{transformer-medium-hand} is the same network with the
hand-crafted features added to its context features.
Table~\ref{tab:settings} in the appendix lists the hyperparameters.
\fi

\paragraph{Training.}
All networks are trained in the same way: with AdamW~\cite{adamw}, a learning
rate that increases linearly over the first \V{networks.warmup.percent}\% of
the optimisation steps and then follows a cosine schedule to zero, gradient
clipping, and mixed precision on a GPU. The loss matches the primary metric:
mean squared error for a target scored by RMSE, a Huber loss for a target
scored by MAE, and binary cross-entropy for a probability. A regression
target is standardised with the mean and standard deviation of the training
set. After every epoch the network is evaluated on the validation set. We
keep the weights of the best epoch, and stop training when
\V{networks.patience} consecutive epochs have not improved on it.
\ifmore
A model trained on a larger training set is given proportionally fewer
epochs, and a correspondingly shorter patience
(Section~\ref{sec:ablations}).
\fi

\begin{table}[htbp]\centering\footnotesize
\block{runs}
\caption{The models, best first. \emph{Selected} is the number of trees or
the training epoch chosen on the validation set. \emph{Time} is the total
wall-clock time for training and evaluation.}
\label{tab:runs}
\end{table}


\section{Evaluation Protocol}
\label{sec:protocol}

\paragraph{Metrics.}
The primary metric, fixed before the first full training run, is the mean
absolute error of the target on the test set, where lower is better. Because
the target is a logarithm, an absolute error of $e$ means that one plus the
predicted number of seconds and one plus the observed number differ by a
factor of $\exp(e)$: an error of 0.69 is a factor of two. We report
$\exp(\mathrm{MAE})$ alongside the MAE and call it the \emph{error factor}.
It is the geometric mean of this factor over the test positions. We report
three secondary metrics: the root mean squared error (RMSE) of the target,
the coefficient of determination R\textsuperscript{2}, which is the share of
the variance of the target that a model explains, and Spearman's rank
correlation between the predicted and the observed time. We also report
results separately for \V{bands.count} bands of the mover's rating.

\paragraph{Times in seconds.}
For readability we give times in seconds. A mean of observed times or of
predictions over a group of positions is computed on the logarithmic scale
and converted back, as $\exp(\bar{y}) - 1$ for a mean $\bar{y}$ of the target
or of the predictions. For observed times we call this value the
\emph{typical time} of the group, and for predictions the \emph{mean
prediction}. The typical time is the geometric mean of one plus the time,
less one. For a distribution with a long right tail it lies close to the
median and below the arithmetic mean, which the long thinks raise\iflook:
over the whole dataset the typical time is \V{label.mean} seconds, the median
is \V{describe.spent.q50} seconds and the arithmetic mean is
\V{describe.spent.mean} seconds\fi.

% Evaluation protocol: the same for every task. Each model paper names its
% own metrics under its section heading and inputs the rest from here.
% The prose follows docs/writing.md.
\paragraph{Confidence intervals.}
Examples from the same game are not independent, so we resample whole games.
The test games are drawn with replacement \V{method.resamples} times, the
metric is computed each time on the examples of the games drawn, and the
2.5th and 97.5th percentiles of these values form the 95\% confidence
interval (CI) that we report with every score.

\paragraph{Paired comparisons.}
To compare two models we compute both scores on the same resampled games and
take their difference. We call this difference the \emph{gain} of the first
model over the second, and sign it so that a positive gain means that the
first model is better. A gain is \emph{statistically significant} when its
95\% CI excludes zero, and is then set in bold in the tables. Unless stated
otherwise, the second model is the GBDT baseline. Two models tend to do well
and badly on the same games, so a paired gain can be significant even when
the two models' separate intervals overlap. These intervals reflect the
sampling of test games only. Each network was trained with a single random
seed, and the variation between training runs is not included.

\paragraph{Model selection.}
\ifchance
The number of trees, the training epoch whose weights are kept and the
calibration parameters are chosen on the validation set.
\else
The number of trees and the training epoch whose weights are kept are chosen
on the validation set.
\fi
The test set is used once per model, for the scores reported here.

\ifchance
\paragraph{Calibration.}
Before a model is scored, its predicted probabilities are rescaled by Platt
scaling~\cite{platt}: a logistic function of the log-odds whose slope and
intercept are fitted on the validation set. The transformation is monotone,
so it leaves a model's ranking of the examples unchanged and affects only its
log loss and its calibration. What remains of a model's miscalibration after
this step cannot be removed by a linear correction of its log-odds.
\fi

\ifpairs
\paragraph{Acceptance criterion.}
Section~\ref{sec:ablations} evaluates modifications of the models. A
modification is \emph{accepted} only if the modified model is better than the
unmodified one on the validation set and its paired gain on the test set is
statistically significant. A modification that is not accepted is still
reported, as a measurement of its effect.
\fi

\paragraph{Tables and figures.}
A model has the same colour and marker in every figure, and the GBDT baseline
is always drawn in orange. The values plotted in a figure are also given in a
table, in the text or in the appendix. All numbers, tables and figures are
generated from the committed result files by the command that builds the
paper.


\section{Results}
\label{sec:results}

\subsection{Main results}

Table~\ref{tab:leaderboard} and Figure~\ref{fig:scores} give the main
results. The constant baseline predicts the median of the training targets,
which corresponds to \V{floor.constant} seconds, for every position, and has
an MAE of \V{floor.score} (an error factor of \V{floor.mae.factor}). GBDT on
the context features alone reaches \V{run.gbm-context.score} and thereby
achieves \V{context.lift.share}\% of the improvement of the GBDT baseline
over the constant baseline. The ratings, the clocks, the time control and the
ply thus carry most of the signal that the models find, before anything is
known about the pieces on the board. The GBDT baseline, with the hand-crafted
features of
the position, reaches \V{reference.score} (an error factor of
\V{reference.mae.factor}).

\begin{table}[htbp]\centering\small
\block{leaderboard}
\caption{Main results on the \V{data.rows.test} test positions, best first.
The gain is the paired difference in MAE from \V{task.reference} on the same
resampled games. It is positive when the model is better, and bold when
statistically significant.}
\label{tab:leaderboard}
\end{table}

\begin{figure}[htbp]\centering
\block{scores}
\caption{Test MAE of every model with its 95\% CI. The vertical line marks
the GBDT baseline.}
\label{fig:scores}
\end{figure}

Networks trained from random initialisation on the board and the context
features are close to the GBDT baseline. Of the \V{ladder.count} networks of
the fixed sequence, \V{ladder.level} are not significantly different from it,
\V{ladder.behind} is significantly worse and \V{ladder.ahead} is
significantly better (Figure~\ref{fig:gains}). The small ResNet (gain
\V{run.conv-small.gain}; 95\% CI \V{run.conv-small.gain.low} to
\V{run.conv-small.gain.high}) and the small transformer
(\V{run.transformer-small.gain}; \V{run.transformer-small.gain.low} to
\V{run.transformer-small.gain.high}) are the two that do not differ
significantly from the baseline. The medium transformer is significantly
worse (\V{run.transformer-medium.gain}; \V{run.transformer-medium.gain.low}
to \V{run.transformer-medium.gain.high}), although it has
\V{run.transformer-medium.parameters.millions} million parameters and the
small models have \V{run.conv-small.parameters.millions} and
\V{run.transformer-small.parameters.millions} million.\ifchangeepochs{}
Section~\ref{sec:ablations} shows that the training schedule is responsible:
with twice the number of epochs the same model is significantly better than
the baseline.\fi

\begin{figure}[htbp]\centering
\block{gains}
\caption{Paired gain in MAE of each network over the GBDT baseline, with its
95\% CI. A model whose interval lies entirely to the right of the line is
significantly better than the baseline, and one whose interval lies entirely
to the left is significantly worse.}
\label{fig:gains}
\end{figure}

The hand-crafted features improve the medium transformer.
\run{transformer-medium-hand}, the medium transformer with the hand-crafted
features added to its context features, is significantly better than the
GBDT baseline (gain \V{run.transformer-medium-hand.gain}; 95\% CI
\V{run.transformer-medium-hand.gain.low} to
\V{run.transformer-medium-hand.gain.high}).
\ifpairs
Compared with \run{transformer-medium}, it is
\V{pair.transformer-medium-hand.transformer-medium.verdict} on the test set
(gain \V{pair.transformer-medium-hand.transformer-medium.gain}; 95\% CI
\V{pair.transformer-medium-hand.transformer-medium.gain.low} to
\V{pair.transformer-medium-hand.transformer-medium.gain.high}).
\fi
A network that reads the board therefore adds information to the hand-crafted
features.

The remaining models in Table~\ref{tab:leaderboard} are modifications that
passed the acceptance criterion of Section~\ref{sec:protocol}, and
Section~\ref{sec:ablations} describes them\ifchangestart: two transformers
with a pretrained encoder, which are the two best models\fi\ifshapes, a
ResNet and a transformer with an attack-aware attention
bias\fi\ifchangeepochs, two transformers with a longer training
schedule\fi\ifmore, and models trained on a larger training set\fi. Compared
with the GBDT baseline, \V{networks.ahead.of} on the leaderboard are
significantly better. The best model's gain is \V{best.gain.percent}\% of the
GBDT baseline's MAE, and \V{best.gain.of-lift}\% of the improvement of the
GBDT baseline over the constant baseline. In terms of the error factor, the
best model reaches \V{best.mae.factor}, compared with
\V{reference.mae.factor} for the GBDT baseline and \V{floor.mae.factor} for
the constant baseline. GBDT on the context and the hand-crafted features
therefore accounts for most of what is predicted, and the best network adds a
smaller improvement.

\subsection{Explained variance and rank correlation}

Table~\ref{tab:metrics} gives all metrics for every model. The
R\textsuperscript{2} of the best model is \V{best.r2}, compared with
\V{reference.r2} for the GBDT baseline and \V{run.gbm-context.r2} for GBDT on
the context features alone. Its rank correlation with the observed time is
\V{best.spearman}, compared with \V{reference.spearman} and
\V{run.gbm-context.spearman}. The RMSE of the best model is \V{best.rmse},
compared with \V{reference.rmse} for the GBDT baseline and \V{floor.rmse} for
the constant baseline. The four metrics order the models in nearly the same
way, and the same two pretrained models are the best by each of them.

\begin{table}[htbp]\centering\small
\block{metrics}
\caption{All metrics for every model on the test set. MAE and RMSE are
errors of the target and lower is better. For R\textsuperscript{2} and
Spearman's rank correlation higher is better. The constant baseline predicts
a single value and has no rank correlation.}
\label{tab:metrics}
\end{table}

Most of the variance of the target is thus not explained by any model, and
we consider this the main finding of the paper. A plausible explanation is
that many reasons for a fast or a slow move lie outside the inputs. A move
may be played at once because it was planned during the opponent's turn or
entered in advance. It may take long because the opponent's move refuted the
player's plan, or because the player was interrupted. The dataset records
none of these circumstances, so we cannot separate them. What a model of a
single position can estimate is the usual time for a kind of position at a
given rating and clock, and Section~\ref{sec:error} shows that the best model
estimates it accurately on average.

\subsection{Results by rating band}

Table~\ref{tab:bands} gives the MAE by band of the mover's rating. The MAE of
the constant baseline rises with the rating, from \V{band.1.floor} for movers
rated \V{band.1.name} to \V{band.5.floor} for movers rated \V{band.5.name},
and the MAE of every other model rises with it. The constant baseline
predicts one value for all positions, so its MAE in a band measures how far
the times of that band lie from this value.\iflook{} The typical time differs
little between the bands, from \V{slice.band.-1100.happened} seconds for
movers rated below 1100 to \V{slice.band.2000+.happened} seconds for movers
rated 2000 and above, so the increase mainly reflects a wider spread: the
times of stronger players vary more.\fi{} A higher MAE in a higher band
therefore does not show that a model is less accurate there relative to what
can be predicted, and scores should be compared between models within a band
and not between bands.

\begin{table}[htbp]\centering\small
\block{bands-table}
\caption{Test MAE by band of the mover's rating. The spread of the times
differs between bands, so values are comparable within a column and not
between columns.}
\label{tab:bands}
\end{table}

The best model has a lower MAE than the GBDT baseline in
\V{bands.best-ahead.of}. The difference is largest for movers rated
\V{bands.lead.widest.name} (\V{bands.lead.widest}) and smallest for movers
rated \V{bands.lead.narrowest.name} (\V{bands.lead.narrowest};
Figure~\ref{fig:bands}).\iflook{} With a paired comparison within each band,
the gain is statistically significant in \V{slice.band.ahead.of}
(Table~\ref{tab:slice-band} in the appendix).\fi

\begin{figure}[htbp]\centering
\block{bands}
\caption{Difference in MAE between each network and the GBDT baseline, by
band of the mover's rating (the differences of the values in
Table~\ref{tab:bands}, without intervals). Above the dashed line the network
is better.}
\label{fig:bands}
\end{figure}

\subsection{Training curves}

Figure~\ref{fig:curves} shows the validation MAE of each network of the fixed
sequence after every epoch, and \V{ladder.at-last-epoch.of} reached their
best validation MAE in the final epoch: \run{conv-small} at
\V{run.conv-small.best-epoch} of \V{run.conv-small.epochs},
\run{transformer-small} at \V{run.transformer-small.best-epoch} of
\V{run.transformer-small.epochs} and \run{transformer-medium} at
\V{run.transformer-medium.best-epoch} of \V{run.transformer-medium.epochs}.
The validation MAE of these models was still decreasing when training ended.
\run{transformer-medium-hand} reached its best epoch at
\V{run.transformer-medium-hand.best-epoch} of
\V{run.transformer-medium-hand.epochs}.
\ifchangeepochs
Section~\ref{sec:ablations} shows that a longer schedule improves both medium
transformers.
\fi

\begin{figure}[htbp]\centering
\block{curves}
\caption{Validation MAE of each network of the fixed sequence after every
epoch. The dashed line is the GBDT baseline's validation MAE.}
\label{fig:curves}
\end{figure}

\subsection{Feature importance of the GBDT baseline}

The GBDT baseline consists of \V{reference.trees} trees over
\V{reference.features} features. Figure~\ref{fig:leans} shows each feature's
share of the total split gain. The three most important features are
\V{reference.leans.1.feature} (\V{reference.leans.1.share}\% of the gain),
\V{reference.leans.2.feature} (\V{reference.leans.2.share}\%) and
\V{reference.leans.3.feature} (\V{reference.leans.3.share}\%). The clock is
the most important feature by a wide margin, and the most important feature
of the position is the number of moves available to the player. Split gain
shows which features the model used. It does not identify what makes a
player think, because correlated features divide the credit between them
arbitrarily.

\begin{figure}[htbp]\centering
\block{leans}
\caption{The ten most important features of the GBDT baseline, by share of
the total split gain.}
\label{fig:leans}
\end{figure}

\ifchanges
\section{Ablations}
\label{sec:ablations}

% Ablations: what a modification is, the table of paired comparisons and the
% two figures of validation curves. The same for every task, so a model paper
% inputs this file at the start of its ablation section and then reports its
% own findings. Each part appears once the task has the runs it needs, and the
% list below needs at least one of its items.
% The prose follows docs/writing.md.
\ifpairs
The main results compare complete models. An ablation isolates one factor: it
compares a model with its \emph{base}, a model that differs from it in a
single respect. Table~\ref{tab:pairs} lists every such comparison for this
task, with the gain of the modified model over its base on the validation set
and its paired gain on the test set, computed on the same resampled test
games as in the main results. The acceptance criterion of
Section~\ref{sec:protocol} is applied to these pairs. The modifications are
the following.
\begin{itemize}
\ifchangehand
\item \textbf{Hand-crafted features} (suffix \code{-hand}): the network
receives the hand-crafted features as additional context features.
\fi
\ifchangeepochs
\item \textbf{Longer schedule} (suffix \code{-long}): twice the training
epochs, with the same warm-up fraction and cosine schedule.
\fi
\ifchangestart
\item \textbf{Pretraining} (suffix \code{-pretrained}): the transformer
encoder is initialised from a transformer of the same size that was trained
to predict the engine's win probability for positions from the Lichess
evaluation database~\cite{lichessdb}. The piece and square embeddings and all
transformer layers are copied. The prediction head is initialised randomly,
as are the projections of the auxiliary planes that the pretraining task does
not have and, unless the two tasks have the same context features, the
projection of the context features. Training then follows the standard
schedule.
\fi
\ifchangeattacks
\item \textbf{Attack-aware attention bias} (suffix \code{-attacks}): the
relative position bias encodes where two squares lie and not what stands
between them. With this modification each attention head learns four
additional biases for a pair of squares, applied when the piece on one square
attacks the other, when it is attacked from the other, and for the same two
relations through exactly one intervening piece, as in a pin or a skewer. The
relations are computed from the tokens without search, and the biases are
initialised to zero.
\fi
\ifchangegeometry
\item \textbf{No relative position bias} (suffix \code{-no-geometry}): the
relative position bias is removed. This configuration measures the
contribution of the bias and is not a candidate model.
\fi
\ifchangeconv
\item \textbf{ResNet encoder} (prefix \code{conv-medium}): a ResNet with
squeeze-and-excitation in place of the transformer, trained at the higher
peak learning rate used for the convolutional models.
\fi
\ifchangehybrid
\item \textbf{Hybrid encoder} (prefix \code{hybrid-medium}): a few residual
convolutional blocks followed by transformer layers over the resulting
square representations.
\fi
\ifchangelarger
\item \textbf{Larger transformer} (prefix \code{transformer-large}): more layers and
a larger width, with the training schedule unchanged.
\fi
\ifgames
\item \textbf{Training-set size}: the same model trained on a fraction or a
multiple of the training games (Section~\ref{sec:training-size}).
\fi
\iffeatures
\item \textbf{Feature groups}: a model trained without one group of
hand-crafted features at a time (Section~\ref{sec:feature-groups}).
\fi
\end{itemize}

\longblock{pairs}{Paired comparisons. Each row compares a model with a
baseline, in most rows the model it modifies. \emph{Change} states how the
two differ. The gain is that of the model over the baseline and is positive
when the model is better. On
the validation set it is the difference between the two scores. On the test
set it is the paired gain, with its 95\% CI, and is set in bold when
statistically significant.}{tab:pairs}
\fi

\ifvariants
Figure~\ref{fig:variants} shows the validation curves of each model of the
fixed sequence whose training schedule or initialisation was modified,
together with the modified models. A modified model has the colour of its
base and a broken line: long dashes for a longer schedule and dots for a
pretrained encoder.

\begin{figure}[htbp]\centering
\block{variants}
\caption{Validation curves of a model of the fixed sequence (solid line) and
of the models that modify its training schedule or initialisation (broken
lines), after every epoch. One panel per base model, on a common scale. The
dashed horizontal line is the GBDT baseline's validation score.}
\label{fig:variants}
\end{figure}
\fi

\ifshapes
Figure~\ref{fig:shapes} shows the validation curves of the models that
change the architecture and nothing about the training data, each in its own
panel beside its base.

\begin{figure}[htbp]\centering
\block{shapes}
\caption{Validation curves of a model with a modified architecture (broken
line) and of its base (solid line), after every epoch. One panel per modified
model, on a common scale. The dashed horizontal line is the GBDT baseline's
validation score.}
\label{fig:shapes}
\end{figure}
\fi


\ifchangehand
\subsection{Hand-crafted features}

The hand-crafted features improve the medium transformer. Compared with
\run{transformer-medium}, \run{transformer-medium-hand} is
\V{pair.transformer-medium-hand.transformer-medium.valid.verdict} than its
base on the validation set and
\V{pair.transformer-medium-hand.transformer-medium.verdict} on the test set
(gain \V{pair.transformer-medium-hand.transformer-medium.gain} MAE; 95\% CI
\V{pair.transformer-medium-hand.transformer-medium.gain.low} to
\V{pair.transformer-medium-hand.transformer-medium.gain.high}), so the
modification is \V{pair.transformer-medium-hand.transformer-medium.kept}.
\ifchangeepochs
The result is the same under the longer schedule of the next subsection:
compared with \run{transformer-medium-long},
\run{transformer-medium-hand-long} is
\V{pair.transformer-medium-hand-long.transformer-medium-long.verdict} on the
test set (gain
\V{pair.transformer-medium-hand-long.transformer-medium-long.gain}; 95\% CI
\V{pair.transformer-medium-hand-long.transformer-medium-long.gain.low} to
\V{pair.transformer-medium-hand-long.transformer-medium-long.gain.high}).
\fi
With the schedule of the fixed sequence, the features move the medium
transformer from significantly below the GBDT baseline to significantly above
it: its gains over the baseline are \V{run.transformer-medium.gain} without
them and \V{run.transformer-medium-hand.gain} with them.
\fi

\ifchangeepochs
\subsection{Training schedule}

Doubling the number of epochs, from \V{run.transformer-medium.epochs} to
\V{run.transformer-medium-long.epochs}, improves both medium transformers.
Without hand-crafted features, \run{transformer-medium-long} is
\V{pair.transformer-medium-long.transformer-medium.valid.verdict} than its
base on the validation set and
\V{pair.transformer-medium-long.transformer-medium.verdict} on the test set
(gain \V{pair.transformer-medium-long.transformer-medium.gain}; 95\% CI
\V{pair.transformer-medium-long.transformer-medium.gain.low} to
\V{pair.transformer-medium-long.transformer-medium.gain.high}), so the
modification is \V{pair.transformer-medium-long.transformer-medium.kept}.
With hand-crafted features, \run{transformer-medium-hand-long} is
\V{pair.transformer-medium-hand-long.transformer-medium-hand.valid.verdict}
than its base on the validation set and
\V{pair.transformer-medium-hand-long.transformer-medium-hand.verdict} on the
test set (gain
\V{pair.transformer-medium-hand-long.transformer-medium-hand.gain}; 95\% CI
\V{pair.transformer-medium-hand-long.transformer-medium-hand.gain.low} to
\V{pair.transformer-medium-hand-long.transformer-medium-hand.gain.high}),
and the modification is
\V{pair.transformer-medium-hand-long.transformer-medium-hand.kept}.

The longer schedule matters more for the model without hand-crafted features:
its paired gain is the larger of the two, and the two intervals do not
overlap. With the longer schedule, the gain of that model over the GBDT
baseline is \V{run.transformer-medium-long.gain} (95\% CI
\V{run.transformer-medium-long.gain.low} to
\V{run.transformer-medium-long.gain.high}), where it was
\V{run.transformer-medium.gain} with the schedule of the fixed sequence. The
result of the fixed sequence, in which the medium transformer without
hand-crafted features is significantly worse than the GBDT baseline,
therefore reflects its schedule: the model was undertrained. Both models with
the longer schedule reached their best validation epoch near its end,
\run{transformer-medium-long} at \V{run.transformer-medium-long.best-epoch}
of \V{run.transformer-medium-long.epochs} and
\run{transformer-medium-hand-long} at
\V{run.transformer-medium-hand-long.best-epoch}
(Figure~\ref{fig:variants}), so a still longer schedule may improve them
further. We have not tried one.
\fi

\ifchangestart
\subsection{Pretraining}

A pretrained encoder improves both medium transformers, and by more than any
other modification. Without hand-crafted features,
\run{transformer-medium-pretrained} is
\V{pair.transformer-medium-pretrained.transformer-medium.valid.verdict} than
its base on the validation set and
\V{pair.transformer-medium-pretrained.transformer-medium.verdict} on the
test set (gain \V{pair.transformer-medium-pretrained.transformer-medium.gain}
MAE; 95\% CI
\V{pair.transformer-medium-pretrained.transformer-medium.gain.low} to
\V{pair.transformer-medium-pretrained.transformer-medium.gain.high}), and
the modification is
\V{pair.transformer-medium-pretrained.transformer-medium.kept}. With
hand-crafted features, \run{transformer-medium-hand-pretrained} is
\V{pair.transformer-medium-hand-pretrained.transformer-medium-hand.valid.verdict}
than its base on the validation set and
\V{pair.transformer-medium-hand-pretrained.transformer-medium-hand.verdict}
on the test set (gain
\V{pair.transformer-medium-hand-pretrained.transformer-medium-hand.gain};
95\% CI
\V{pair.transformer-medium-hand-pretrained.transformer-medium-hand.gain.low}
to
\V{pair.transformer-medium-hand-pretrained.transformer-medium-hand.gain.high}),
and the modification is
\V{pair.transformer-medium-hand-pretrained.transformer-medium-hand.kept}.
The best model of this task is \V{best.name}.

With a pretrained encoder the hand-crafted features contribute little. The
pretrained model with hand-crafted features has a test MAE of
\V{run.transformer-medium-hand-pretrained.score} and the pretrained model
without them one of \V{run.transformer-medium-pretrained.score}, a
difference of \V{run.transformer-medium-hand-pretrained.over-twin}, where
the paired gain from the features is
\V{pair.transformer-medium-hand.transformer-medium.gain} from random
initialisation. The two pretrained models were not compared directly on
resampled games, and their gains over the GBDT baseline lie within each
other's intervals (Table~\ref{tab:leaderboard}). We hypothesise that an
encoder pretrained to predict engine evaluations already represents most of
what the hand-crafted features provide.

The pretrained model with hand-crafted features reached its best validation
epoch at \V{run.transformer-medium-hand-pretrained.best-epoch} of
\V{run.transformer-medium-hand-pretrained.epochs}, and the pretrained model
without them at \V{run.transformer-medium-pretrained.best-epoch}
(Figure~\ref{fig:variants}). The validation MAE of neither model was still
decreasing when its schedule ended.
\fi

\ifshapes
\subsection{Architecture}

The architecture ablations keep the training data, the prediction head and
the number of epochs of a medium transformer and change the encoder. All of
these models are trained from random initialisation, as their bases are.
\begin{itemize}
\ifchangeconv
\item \textbf{ResNet.} \run{conv-medium-hand} has
\V{run.conv-medium-hand.parameters.millions} million parameters, compared
with \V{run.transformer-medium-hand.parameters.millions} million for the
transformer, and is trained with a peak learning rate of
\V{run.conv-medium-hand.lr} instead of \V{run.transformer-medium-hand.lr}.
Compared with \run{transformer-medium-hand} it is
\V{pair.conv-medium-hand.transformer-medium-hand.valid.verdict} than its base
on the validation set and
\V{pair.conv-medium-hand.transformer-medium-hand.verdict} on the test set
(gain \V{pair.conv-medium-hand.transformer-medium-hand.gain}; 95\% CI
\V{pair.conv-medium-hand.transformer-medium-hand.gain.low} to
\V{pair.conv-medium-hand.transformer-medium-hand.gain.high}), and the
modification is \V{pair.conv-medium-hand.transformer-medium-hand.kept}.
\fi
\ifchangehybrid
\item \textbf{Hybrid.} \run{hybrid-medium-hand} has
\V{run.hybrid-medium-hand.convs} residual convolutional blocks followed by
\V{run.hybrid-medium-hand.depth} transformer layers, for a total of
\V{run.hybrid-medium-hand.parameters.millions} million parameters. Compared
with \run{transformer-medium-hand} it is
\V{pair.hybrid-medium-hand.transformer-medium-hand.verdict} on the test set
(gain \V{pair.hybrid-medium-hand.transformer-medium-hand.gain}; 95\% CI
\V{pair.hybrid-medium-hand.transformer-medium-hand.gain.low} to
\V{pair.hybrid-medium-hand.transformer-medium-hand.gain.high}), and the
modification is \V{pair.hybrid-medium-hand.transformer-medium-hand.kept}.
\fi
\ifchangeattacks
\item \textbf{Attack-aware attention bias.} Compared with
\run{transformer-medium}, which has no hand-crafted features,
\run{transformer-medium-attacks} is
\V{pair.transformer-medium-attacks.transformer-medium.valid.verdict} than its
base on the validation set and
\V{pair.transformer-medium-attacks.transformer-medium.verdict} on the test
set (gain \V{pair.transformer-medium-attacks.transformer-medium.gain}; 95\% CI
\V{pair.transformer-medium-attacks.transformer-medium.gain.low} to
\V{pair.transformer-medium-attacks.transformer-medium.gain.high}), and the
modification is
\V{pair.transformer-medium-attacks.transformer-medium.kept}. Compared with
\run{transformer-medium-hand}, \run{transformer-medium-hand-attacks} is
\V{pair.transformer-medium-hand-attacks.transformer-medium-hand.valid.verdict}
than its base on the validation set and
\V{pair.transformer-medium-hand-attacks.transformer-medium-hand.verdict} on
the test set (gain
\V{pair.transformer-medium-hand-attacks.transformer-medium-hand.gain}; 95\%
CI \V{pair.transformer-medium-hand-attacks.transformer-medium-hand.gain.low}
to
\V{pair.transformer-medium-hand-attacks.transformer-medium-hand.gain.high}),
and the modification is
\V{pair.transformer-medium-hand-attacks.transformer-medium-hand.kept},
because the criterion requires an improvement on both sets.
\fi
\ifchangegeometry
\item \textbf{No relative position bias.} Compared with
\run{transformer-medium-hand},
\run{transformer-medium-hand-no-geometry} is
\V{pair.transformer-medium-hand-no-geometry.transformer-medium-hand.verdict}
on the test set (gain
\V{pair.transformer-medium-hand-no-geometry.transformer-medium-hand.gain};
95\% CI
\V{pair.transformer-medium-hand-no-geometry.transformer-medium-hand.gain.low}
to
\V{pair.transformer-medium-hand-no-geometry.transformer-medium-hand.gain.high}).
Removing the bias has no measurable effect on this task.
\fi
\end{itemize}

\ifchangeattacks
Among the modifications of the medium transformer that keep its random
initialisation, the attack-aware bias gives the largest paired gain
(\V{pair.transformer-medium-attacks.transformer-medium.gain}), more than the
hand-crafted features do
(\V{pair.transformer-medium-hand.transformer-medium.gain}). The test MAE of
\run{transformer-medium-attacks} is \V{run.transformer-medium-attacks.score},
compared with \V{run.transformer-medium-hand.score} for the transformer with
the features. The two models were not compared directly on resampled games,
but the intervals of their gains over the GBDT baseline do not overlap
(Table~\ref{tab:leaderboard}; \V{run.transformer-medium-attacks.gain.low} to
\V{run.transformer-medium-attacks.gain.high}, and
\V{run.transformer-medium-hand.gain.low} to
\V{run.transformer-medium-hand.gain.high}). The two additions do not
combine. With both, the test MAE is
\V{run.transformer-medium-hand-attacks.score}, which is higher than with the
bias alone, and the interval of this model's gain over the GBDT baseline
(\V{run.transformer-medium-hand-attacks.gain.low} to
\V{run.transformer-medium-hand-attacks.gain.high}) lies below that of the
model with the bias alone.
\fi

\ifchangeconv
The ResNet is the best model trained from random initialisation on the
standard training set (\V{run.conv-medium-hand.score}). Its comparison with
the transformer is not a comparison of the encoders alone, because the two
models are trained with different learning rates.
\fi

We do not know why these two encoders are better on this task. One possible
explanation is that the time a player spends depends on local tactical
content, such as the captures that are available, and that a convolution over
neighbouring squares or a bias along the lines of attack makes this content
easier to learn than unconstrained attention does. The error analysis is
consistent with the first part of this explanation\iflook: the gain of the
best network over the GBDT baseline increases with the number of legal
captures (Section~\ref{sec:error})\fi. It does not test the second part. The
same two modifications are accepted for the puzzle-difficulty task of this
repository, and neither differs significantly from its base for the
blunder-risk task, which uses the positions of this paper. The explanation is
therefore incomplete at best. Each of these models was trained once.
\fi

% Training-set size: the same ablation for every task, so a model paper
% inputs this file inside its ablation section and then reports its own
% findings. It appears once the task has the runs.
% The prose follows docs/writing.md.
\ifgames
\subsection{Training-set size}
\label{sec:training-size}

To estimate how much a model's accuracy depends on the amount of training
data, we retrain it on a half and on a quarter of the training games and
evaluate it on the full test set (Table~\ref{tab:games},
Figure~\ref{fig:games}). The subsets are selected by a hash of the game
identifier and are nested, so the quarter is contained in the half. A network
is given twice the epochs on half the games and four times the epochs on a
quarter, so the number of optimisation steps is constant and only the amount
of data varies.
\ifmore
We also train on more data. A larger sample of the same source is split in
the same way as the standard dataset, so the validation and test sets are
unchanged and no test game appears in the larger training set. A network
trained on a larger set is given proportionally fewer epochs, so that it
again takes the same number of optimisation steps as its base. Unlike a model
trained on a subset, a model trained on more data is a candidate model, and
the acceptance criterion applies to it.
\fi

\begin{table}[htbp]\centering\footnotesize
\block{study-games}
\caption{Training-set size. Each model is trained on a fraction of the
training games\ifmore, or on a multiple of them from the larger sample,\fi{}
and evaluated on the full test set. The gain is paired with the same model
trained on the standard training set. A negative gain means that the model is
worse, and bold marks statistical significance.}
\label{tab:games}
\end{table}

\begin{figure}[htbp]\centering
\block{study-games-figure}
\caption{Test score as a function of the number of training games. Each step
to the right doubles the training set.}
\label{fig:games}
\end{figure}
\fi


\ifgames
\paragraph{Smaller training sets.}
Both model classes are less accurate with less data, and the transformer
loses more. Relative to the same model trained on the full training set,
\run{transformer-medium-hand} has a paired gain of
\V{pair.transformer-medium-hand-half.transformer-medium-hand.gain} MAE on
half of the training games and of
\V{pair.transformer-medium-hand-quarter.transformer-medium-hand.gain} on a
quarter, and it is significantly worse in
\V{study.games.transformer-medium-hand.worse.of}. For GBDT the gains are
\V{pair.gbm-hand-half.gbm-hand.gain}
(\V{pair.gbm-hand-half.gbm-hand.verdict}) and
\V{pair.gbm-hand-quarter.gbm-hand.gain}
(\V{pair.gbm-hand-quarter.gbm-hand.verdict}). On a quarter of the games the
transformer still has the lower MAE of the two
(\V{run.transformer-medium-hand-quarter.score}, compared with
\V{run.gbm-hand-quarter.score} for GBDT), and its advantage over GBDT grows
with each doubling up to the full training set (Table~\ref{tab:games},
Figure~\ref{fig:games}).
\fi

\ifmore
\paragraph{Larger training sets.}
The larger sample contains up to \V{more.gbm-hand.most} the training games,
with the validation and test sets unchanged. Relative to the same model on
the standard training set, \run{transformer-medium-hand} has a paired gain of
\V{pair.transformer-medium-hand-double.transformer-medium-hand.gain} MAE on
twice the training games
(\V{pair.transformer-medium-hand-double.transformer-medium-hand.verdict};
\V{pair.transformer-medium-hand-double.transformer-medium-hand.kept}) and of
\V{pair.transformer-medium-hand-quadruple.transformer-medium-hand.gain} on
four times
(\V{pair.transformer-medium-hand-quadruple.transformer-medium-hand.verdict};
\V{pair.transformer-medium-hand-quadruple.transformer-medium-hand.kept}).
For GBDT the gains are \V{pair.gbm-hand-double.gbm-hand.gain}
(\V{pair.gbm-hand-double.gbm-hand.verdict};
\V{pair.gbm-hand-double.gbm-hand.kept}) and
\V{pair.gbm-hand-quadruple.gbm-hand.gain}
(\V{pair.gbm-hand-quadruple.gbm-hand.verdict};
\V{pair.gbm-hand-quadruple.gbm-hand.kept}).

More data helps both model classes, and it helps GBDT more: the gain of GBDT
is the larger one at both sizes, and on four times the games the intervals of
the two gains do not overlap (Table~\ref{tab:games}). The transformer's
advantage over GBDT, which grows up to the standard training set, is
therefore slightly smaller at both larger sizes. With four times the games,
GBDT reaches an MAE of \V{run.gbm-hand-quadruple.score}, which is still
higher than the \V{run.transformer-medium-hand.score} of the transformer on
the standard training set. \run{transformer-medium-hand-quadruple} has an MAE
of \V{run.transformer-medium-hand-quadruple.score}.
\ifchangeconv
On this task the choice of encoder matters more than this amount of data.
Compared with \run{transformer-medium-hand}, the paired gain of the ResNet on
the standard training set
(\V{pair.conv-medium-hand.transformer-medium-hand.gain}) is larger than that
of the transformer on four times the games, and the two intervals do not
overlap (Table~\ref{tab:pairs}).
\fi

The three transformer models take the same number of optimisation steps. On
the standard training set the transformer reached its best validation epoch
at \V{run.transformer-medium-hand.best-epoch} of
\V{run.transformer-medium-hand.epochs}. On the larger training sets it
reached it in the final epoch, at
\V{run.transformer-medium-hand-double.best-epoch} of
\V{run.transformer-medium-hand-double.epochs} on twice the games and at
\V{run.transformer-medium-hand-quadruple.best-epoch} of
\V{run.transformer-medium-hand-quadruple.epochs} on four times, so its
validation MAE was still decreasing when its schedule ended.
\ifchangeepochs
On the standard training set, doubling the number of epochs alone gained
\V{pair.transformer-medium-hand-long.transformer-medium-hand.gain}, which is
of the same order as the gains from the larger training sets.
\fi
We hypothesise that with more data the transformer is limited by its number
of optimisation steps, in which case its gains from the larger training sets
are lower bounds. GBDT was not limited in this way: on four times the games
it used \V{run.gbm-hand-quadruple.trees} of the \V{method.gbm.rounds}
boosting rounds it was allowed.
\fi

% Feature-group ablation: the same ablation for every task, so a model paper
% inputs this file inside its ablation section and then reports its own
% findings. It appears once the task has the runs.
% The prose follows docs/writing.md.
\iffeatures
\subsection{Feature groups}
\label{sec:feature-groups}

\ifsequence
The hand-crafted features fall into three groups: move accuracy, clock use
and the number of available moves.
\else
The hand-crafted features fall into groups: material, mobility, hanging
pieces, passed pawns and king safety, and, where the task gives a line of
moves, the opponent's last move, the first move of the line and the line as a
whole.
\fi
We retrain a model with one group removed at a time
(Table~\ref{tab:features}, Figure~\ref{fig:features}). A significant loss
shows that the model used the group and that no other input provides the same
information. The absence of a loss has two possible explanations: the group
is uninformative, or the same information is available from another group or
from the other inputs. The ablation therefore measures the marginal
contribution of a group given everything else the model receives, which can
be much smaller than the information the group carries on its own.

\begin{table}[htbp]\centering\footnotesize
\block{study-features}
\caption{Feature-group ablation. Each model is trained on all hand-crafted
features except one group. The gain is paired with the model trained on all
features. A negative gain means that the model is worse without the group,
and bold marks statistical significance. Groups are ordered by the size of
the loss.}
\label{tab:features}
\end{table}

\begin{figure}[htbp]\centering
\block{study-features-figure}
\caption{Paired gain of the GBDT baseline without one feature group over the
GBDT baseline with all features, with 95\% CIs. Left of the vertical line,
the model is worse without the group.}
\label{fig:features}
\end{figure}
\fi


\iffeatures
\paragraph{Material, hanging pieces and mobility matter most.}
The ablation was run for the GBDT baseline. It is significantly worse in
\V{study.features.gbm-hand.worse.of}. The largest loss follows the removal of
the \V{study.features.gbm-hand.most.group} group, which consists of the total
material, the material balance and the number of minor and major pieces
(paired gain \V{study.features.gbm-hand.most.gain}). Without the
hanging-piece features the paired gain is
\V{pair.gbm-hand-no-hanging.gbm-hand.gain}, without the mobility features,
which include the numbers of legal moves and of legal captures, it is
\V{pair.gbm-hand-no-mobility.gbm-hand.gain}, and without the two features
that count attacks near the kings it is
\V{pair.gbm-hand-no-kings.gbm-hand.gain}. Without the passed-pawn features
the paired gain is \V{pair.gbm-hand-no-pawns.gbm-hand.gain}
(\V{pair.gbm-hand-no-pawns.gbm-hand.verdict}). No single group accounts for
the contribution of the hand-crafted features. Removing all of them gives
\run{gbm-context}, whose gain relative to the GBDT baseline is
\V{run.gbm-context.gain} (Table~\ref{tab:leaderboard}), several times the
largest loss from a single group.
\fi
\fi

\section{Error Analysis}
\label{sec:error}

% Error analysis: what the analysis is and how it is computed, the same for
% every task. A model paper inputs this file at the start of its error
% analysis, before the subgroups that are specific to its task.
% The prose follows docs/writing.md.
\iflook
The main results identify the best model. They do not show where a model is
accurate and where it is not, which determines how it can be used. We
therefore analyse the test predictions of every model by subgroup. All
results in this section are computed from the \V{look.test-rows} test
examples (analysis of \V{look.date}).

\paragraph{Subgroups.}
We partition the test set by one variable at a time. For each subgroup we
report the number of test examples, the mean observed target, the mean
prediction of a model and the model's score. Two models are compared:
\V{challenger.name}, the best model other than the GBDT baseline, and the
GBDT baseline itself. The gain of the first over the second is computed
within each subgroup as for the whole test set, as a paired difference over
resampled games (\V{look.resamples} resamples), and is omitted for subgroups
with fewer than \V{look.smallest-compared} test examples. A subgroup whose
interval contains zero provides no evidence in either direction. A numeric
variable is divided into ranges, and a range written $a$-$b$ includes $a$ and
excludes $b$.

\paragraph{Two-way tables.}
For a pair of variables we show the mean observed target and the mean
prediction in each cell, on a common colour scale. Cells with fewer than
\V{grid.smallest-cell} test examples are left empty.

\paragraph{Examples.}
Example \ifsequence games\else positions\fi{} are selected by fixed rules and
not by inspection: typical examples chosen by a hash of their key, the
examples with the best model's most extreme predictions in each direction,
and its largest errors in each direction, \V{examples.count} of each. They
illustrate the quantities under discussion and are not evidence of how often
a pattern occurs.
\else
\emph{This version of the paper was built before the error analysis was run.
The analysis partitions the test set by rating, clock and properties of the
position, measures calibration, and shows example positions. It is produced
by \texttt{lab analyse} on the machine that holds the dataset and the test
predictions of every model, and this section appears once its output is
committed.}
\fi


\iflook
\subsection{Distribution of the errors}

The median error factor of the best model is \V{best.close.median}, and 90\%
of its predictions are within a factor of \V{best.close.p90} of the time
spent. \V{best.close.within.1.5}\% of the predictions are within a factor of
1.5, \V{best.close.within.2}\% within a factor of two and
\V{best.close.within.3}\% within a factor of three
(Table~\ref{tab:closeness}). As in the definition of the error factor, these
factors relate one plus the predicted number of seconds to one plus the
observed number. For the GBDT baseline the three shares are
\V{reference.close.within.1.5}\%, \V{reference.close.within.2}\% and
\V{reference.close.within.3}\%.

\begin{table}[htbp]\centering\small
\block{closeness}
\caption{Distribution of the errors on the test set: the median error factor,
the 90th percentile, and the share of predictions within a factor of 1.5, two
and three of the time spent.}
\label{tab:closeness}
\end{table}

The last share limits how a prediction can be used. The remaining moves
differ from the prediction by more than a factor of three in one direction or
the other, and they are too frequent for each of them to be regarded as a
mistake in the use of the clock. All networks and the GBDT baseline have a
similar distribution of errors, which indicates that the dispersion is a
property of the moves and not of a particular model.

\subsection{Calibration}

The best model's predictions are accurate on average over their whole range.
In Figure~\ref{fig:reliability} the test positions are divided into
\V{best.reliability.groups} bins of equal size by the model's prediction. In
the lowest bin the mean prediction is \V{best.reliability.lowest.said}
seconds and the typical time spent is \V{best.reliability.lowest.happened}
seconds. In the highest bin they are \V{best.reliability.highest.said} and
\V{best.reliability.highest.happened} seconds. The largest difference in any
bin is \V{best.reliability.widest-gap} seconds, and it occurs in the highest
bin. The deviations have a consistent direction: in the lower bins the mean
prediction is below the typical time and in the upper bins it is above it, so
the predictions are somewhat more spread out than the typical times. The GBDT
baseline has a narrower
range of predictions, from \V{reference.reliability.lowest.said} to
\V{reference.reliability.highest.said} seconds, and a smaller largest
difference, \V{reference.reliability.widest-gap} seconds.

\begin{figure}[htbp]\centering
\block{reliability}
\caption{Typical time spent against mean predicted time, in seconds. The test
positions are divided into bins of equal size by the prediction, and each
point is one bin. The shaded band is the central 80\% of the times spent in
each bin of the first model. Points on the diagonal are unbiased.}
\label{fig:reliability}
\end{figure}

Predictions for individual moves are much less precise, and the range of the
predictions is narrow. For a given prediction in the middle of the range, the
central 80\% of the times spent span about \V{best.reliability.spread}
seconds (shaded band in Figure~\ref{fig:reliability}). The highest bin, which
contains the twentieth of the test positions with the longest predictions,
has a mean prediction of \V{best.reliability.highest.said} seconds, whereas
one move in twenty takes more than \V{describe.spent.q95} seconds and one in
a hundred more than \V{describe.spent.q99}. The model does not predict that a
long think is about to occur.

Over the whole test set, the mean prediction of the best model corresponds to
\V{overall.transformer-medium-hand-pretrained.said} seconds and that of the
GBDT baseline to \V{overall.gbm-hand.said} seconds, compared with a typical
time of \V{overall.happened} seconds. Both model classes are trained with a
loss that grows linearly with large errors, and such a loss draws a
prediction towards the conditional median and away from the conditional mean.
A plausible explanation of the difference is therefore that, for a given
input, the distribution of the target is skewed towards long times. We have
not tested this.

Conditioning on the time spent instead of the prediction shows regression
towards the mean (Table~\ref{tab:seconds}, Figure~\ref{fig:seconds-bias}).
For moves that took less than two seconds, the mean prediction of
\V{challenger.name} is \V{slice.seconds.-2.challenger.said} seconds. For
moves that took 80 seconds or more, whose typical time is
\V{slice.seconds.80+.happened} seconds, it is
\V{slice.seconds.80+.challenger.said} seconds. The corresponding predictions
of the GBDT baseline are \V{slice.seconds.-2.reference.said} and
\V{slice.seconds.80+.reference.said} seconds. Every model with prediction
error shows this effect, because moves that it cannot distinguish receive
predictions near their common mean. The subgroups of this table are defined
by the target itself, so the table measures the part of the time that could
not be predicted. It does not measure a bias that could be corrected.

\begin{table}[htbp]\centering\footnotesize
\block{slice-seconds}
\caption{Test-set results by the time spent on the move, in seconds. For each
group: the number of test positions, the typical time spent, the mean
prediction and MAE of each of the two models, and the paired gain of the
first model over the second.}
\label{tab:seconds}
\end{table}

\begin{figure}[htbp]\centering
\block{slice-seconds-bias}
\caption{Mean prediction minus typical time spent, in seconds, by the time
spent on the move, for the two models compared. Above the line a model
predicts a longer time than was spent, and below it a shorter one.}
\label{fig:seconds-bias}
\end{figure}

The same table shows where the network and the GBDT baseline differ.
\V{challenger.name} is significantly better in \V{slice.seconds.ahead.of}
(\V{slice.seconds.ahead.groups}), which are the fastest and the slowest
moves, and significantly worse in \V{slice.seconds.behind.of}
(\V{slice.seconds.behind.groups}; gain \V{slice.seconds.5-10.gain}), in the
middle of the distribution. This agrees with the wider range of its
predictions: the network moves further from the usual time than the GBDT
baseline does, which brings it closer to the fastest and the slowest moves
and costs accuracy on ordinary ones.

\subsection{Clock}

The time remaining on the mover's clock is the most important feature of the
GBDT baseline, and Table~\ref{tab:clock} shows its relation to the time
spent. With eight minutes or more remaining the typical time is
\V{slice.clock.480+.happened} seconds, and with five to eight minutes it is
\V{slice.clock.300-480.happened} seconds. Below five minutes it decreases
with the clock: \V{slice.clock.60-120.happened} seconds with one to two
minutes remaining, \V{slice.clock.30-60.happened} seconds with 30 to 60
seconds, and \V{slice.clock.-30.happened} seconds with less than 30 seconds.
The typical time in the first group is lower than in the second. A plausible
explanation is that the first group contains the moves that directly follow
the opening, which are played quickly (Table~\ref{tab:slice-ply}).

\begin{table}[htbp]\centering\footnotesize
\block{slice-clock}
\caption{Test-set results by the time remaining on the mover's clock, in
seconds.}
\label{tab:clock}
\end{table}

Both models follow the typical time from group to group
(Figure~\ref{fig:clock}). The mean prediction of \V{challenger.name} differs
from the typical time by at most \V{slice.clock.challenger.widest-off}
seconds in any group, and that of the GBDT baseline by at most
\V{slice.clock.reference.widest-off}. \V{challenger.name} is significantly
better than the GBDT baseline in \V{slice.clock.ahead.of}
(\V{slice.clock.ahead.groups}). In the two groups between 30 seconds and two
minutes the intervals contain zero. These groups contain
\V{slice.clock.30-60.rows} and \V{slice.clock.60-120.rows} of the
\V{look.test-rows} test positions, because most of a rapid game is played
with ample time.

\begin{figure}[htbp]\centering
\block{slice-clock-said}
\caption{Typical time spent and mean predicted time of the two models
compared, by the time remaining on the mover's clock.}
\label{fig:clock}
\end{figure}

\subsection{Phase of the game}

The time spent follows the course of a game. The typical time is
\V{slice.ply.-20.happened} seconds before ply 20, rises to
\V{slice.ply.30-40.happened} seconds between plies 30 and 40, and falls to
\V{slice.ply.80+.happened} seconds from ply 80 on
(Table~\ref{tab:slice-ply} in the appendix). The number of minor and major
pieces on the board shows the same pattern: the typical time is
\V{slice.pieces.9-12.happened} seconds with nine to twelve such pieces and
\V{slice.pieces.-3.happened} seconds with fewer than three
(Table~\ref{tab:slice-pieces}). Players thus think longest per move in the
middlegame. Two possible reasons for the fast late moves cannot be separated
in these tables: fewer pieces remain to be considered, and less time remains
on the clock.

\V{challenger.name} is significantly better than the GBDT baseline in
\V{slice.ply.ahead.of} by ply and in \V{slice.pieces.ahead.of} by material.
Its gain is smallest late in the game: \V{slice.ply.80+.gain} from ply 80 on
and \V{slice.pieces.-3.gain} with fewer than three minor and major pieces.

\subsection{Number of available moves and captures}

The time spent increases with the number of choices that the position
offers. The typical time is \V{slice.moves.-5.happened}
seconds with fewer than five legal moves and \V{slice.moves.5-15.happened}
seconds with five to fifteen, and it increases with every further group, to
\V{slice.moves.45+.happened} seconds with 45 or more
(Table~\ref{tab:slice-moves}). By the number of legal captures it increases
with every group, from \V{slice.captures.0.happened} seconds with none to
\V{slice.captures.4+.happened} seconds with four or more
(Table~\ref{tab:captures}). A mover in check takes
\V{slice.check.yes.happened} seconds, compared with
\V{slice.check.no.happened} seconds for a mover who is not in check
(Table~\ref{tab:slice-check}), and a check leaves few legal moves.

\begin{table}[htbp]\centering\footnotesize
\block{slice-captures}
\caption{Test-set results by the number of legal captures.}
\label{tab:captures}
\end{table}

The number of the mover's hanging pieces, that is pieces attacked by a less
valuable piece or attacked and not defended, is related to the time spent
less strongly: the typical time is \V{slice.hanging.0.happened} seconds with
none and \V{slice.hanging.3+.happened} seconds with three or more
(Table~\ref{tab:slice-hanging}). In these tables a choice between several
captures is associated with a larger increase in the time spent than a threat
against the mover's own pieces is.

The network's advantage over the GBDT baseline depends on the number of
captures. The gain of \V{challenger.name} is \V{slice.captures.0.gain} in
positions without a legal capture and \V{slice.captures.4+.gain} in positions
with four or more, and it increases with every group in between. By the
number of legal moves the gain is smallest in the group
\V{slice.moves.gain.least.group} (\V{slice.moves.gain.least}), where its
interval contains zero, and it is significant in the other groups
(\V{slice.moves.ahead.groups}). The GBDT baseline receives the numbers of
legal moves and of legal captures as features. The network's additional gain
in positions with many captures therefore indicates that the squares
involved carry information beyond these counts.

\subsection{Win probability before the move}

No model in this paper receives an evaluation, so this analysis shows how a
model that does not know the evaluation behaves across it
(Table~\ref{tab:slice-win}). The typical time is longest in positions that
are unbalanced and not yet decided: \V{slice.win.10-20.happened} and
\V{slice.win.20-40.happened} seconds at win probabilities of 10\% to 20\% and
20\% to 40\% for the mover, and \V{slice.win.60-80.happened} and
\V{slice.win.80-90.happened} seconds at 60\% to 80\% and 80\% to 90\%. It is
shortest at the extremes, \V{slice.win.-10.happened} seconds below 10\% and
\V{slice.win.90+.happened} seconds at 90\% and above. In balanced positions,
between 40\% and 60\%, it is \V{slice.win.40-60.happened} seconds. A
plausible explanation of this lower value is that the balanced group contains
most of the positions from the early part of the game, which are played
quickly, so this table does not separate the evaluation from the phase of
the game.

The mean prediction of \V{challenger.name} minus the typical time is between
\V{slice.win.-10.challenger.off} and \V{slice.win.40-60.challenger.off}
seconds in the four groups below 60\%, and it rises to
\V{slice.win.90+.challenger.off} seconds at 90\% and above, compared with
\V{overall.transformer-medium-hand-pretrained.off} seconds over the whole
test set. For the GBDT baseline the difference is
\V{slice.win.-10.reference.off} seconds below 10\% and
\V{slice.win.90+.reference.off} seconds at 90\% and above. Relative to the
prediction of a model that does not know the evaluation, players thus take
slightly longer when they are worse or level, and slightly less when they are
clearly better. All of these differences are below one second.
\V{challenger.name} is significantly better than the GBDT baseline in
\V{slice.win.ahead.of}.

\subsection{Rating and clock}

Figure~\ref{fig:grid} shows the typical time by the mover's rating and the
time remaining on the mover's clock, together with the mean predictions of
the best model. With five minutes or more remaining, the typical time is
\V{grid.rating-clock.-1100.300+.happened} seconds for movers rated below 1100
and \V{grid.rating-clock.2000+.300+.happened} seconds for movers rated 2000
and above. With less than 30 seconds remaining it is
\V{grid.rating-clock.-1100.-30.happened} and
\V{grid.rating-clock.2000+.-30.happened} seconds. The strongest players thus
speed up most when little time remains, but the two cells contain only
\V{grid.rating-clock.-1100.-30.rows} and \V{grid.rating-clock.2000+.-30.rows}
test positions, and the typical time does not decrease steadily with the
rating in the cells between them. Both models reproduce the joint pattern:
the largest deviation of the mean prediction from the typical time in any
cell is \V{grid.rating-clock.transformer-medium-hand-pretrained.widest-off}
seconds for the best model and \V{grid.rating-clock.gbm-hand.widest-off}
seconds for the GBDT baseline.

\begin{figure}[htbp]\centering
\block{grid-rating-clock}

\medskip
\block{grid-rating-clock-said}
\caption{Typical time spent in seconds by the mover's rating and the time
remaining on the mover's clock. Top: observed times. Bottom: mean predictions
of the best model other than the GBDT baseline. Both panels use the same
colour scale, and cells with too few test positions are left empty.}
\label{fig:grid}
\end{figure}

\subsection{Agreement with the GBDT baseline}

Table~\ref{tab:beside} gives, for each model, the correlation of its
predictions with those of the GBDT baseline on the test set. For the best
model the correlation is \V{best.beside.correlation}, and its prediction is
closer to the time spent than the baseline's on \V{best.beside.closer}\% of
the test positions. The two models therefore make similar predictions, and
the network's gain is a small improvement on a small majority of the moves.

\begin{table}[htbp]\centering\small
\block{beside}
\caption{Agreement with the GBDT baseline on the test set: the correlation
between the two models' predictions, and the share of positions on which the
model's prediction is closer to the time spent than the baseline's.}
\label{tab:beside}
\end{table}

\subsection{Example positions}
\label{sec:positions}

The following examples were selected from the test set by the fixed rules
described above. Each board is shown from the mover's side. The grey arrow is
the opponent's preceding move and the blue arrow is the move that was played.
Next to each board we give the move, the time spent on it, the predictions of
the best model and of the GBDT baseline, and the context of the position. The
context includes the mover's win probability before the move, which no model
receives. The games are public Lichess games.

The largest errors show what the inputs do not contain. The three moves with
the largest overestimates (Figure~\ref{fig:over}) took no measurable time: a
time spent of zero seconds was recorded for each of them. On Lichess a player
can enter a move during the opponent's turn, and such a move is played at
once, which would explain these cases. The three moves with the largest
underestimates (Figure~\ref{fig:under}) each took more than two minutes early
in the game, and two of them are consecutive moves of one player in the same
game. Nothing in the position indicates either kind of event, and a model of
the usual time should not be expected to predict them. The positions with the
highest and the lowest predictions are shown in the appendix
(Figures~\ref{fig:highest} and~\ref{fig:lowest}).

\examples{examples-typical}{Typical positions: the first \V{examples.count}
test positions in the order of a hash of their key.}{fig:typical}

\examples{examples-over}{The largest overestimates: the test positions for
which the best model's prediction exceeds the time spent by the
most.}{fig:over}

\examples{examples-under}{The largest underestimates: the test positions for
which the best model's prediction falls below the time spent by the
most.}{fig:under}
\fi

\section{Practical Use}
\label{sec:use}

\begin{itemize}
\iflook
\item \textbf{A reference time in seconds.} The prediction for a position, a
rating and a clock is accurate on average over its whole range: within
\V{best.reliability.widest-gap} seconds of the typical time in every bin. The
application can show it next to the time that the player spent.
\item \textbf{Statistics over many moves, not a judgement of one move.}
\V{best.close.within.3}\% of the predictions are within a factor of three of
the time spent, so a single move that differs widely from the reference is
ordinary. A review can rely on counts over a game or a month, for example how
often a player moved in less than a third of the reference time and in which
kind of position.
\else
\item \textbf{Statistics over many moves, not a judgement of one move.} The
error factor of the best model is \V{best.mae.factor}, so a single move that
differs widely from the reference is ordinary. A review can rely on counts
over a game or a month.
\fi
\item \textbf{The context provides most of the accuracy.} GBDT on the scalar
context features alone achieves \V{context.lift.share}\% of the improvement
of the GBDT baseline over the constant baseline. A reference computed from
the ratings, the clocks, the time control and the ply already captures most
of what any of the models predicts.
\item \textbf{GBDT provides nearly all of the accuracy.} The GBDT baseline
has an error factor of \V{reference.mae.factor}, compared with
\V{best.mae.factor} for the best model, and needs no GPU. Where the
application has to run the model itself, GBDT is the natural starting point.
\item \textbf{Not a predictor of long thinks.} The model predicts the usual
time for a position and does not identify the individual moves on which a
player thinks for a minute or more\iflook: for the moves that took 80 seconds
or more, its mean prediction is \V{slice.seconds.80+.challenger.said}
seconds\fi. It describes what players do and not what they should do, so it
cannot tell a player where to slow down.
\end{itemize}

\section{Limitations}
\label{sec:limits}

\begin{itemize}
\item \textbf{Much of the variation lies outside the inputs.} A move entered
in advance, the time the opponent has just spent and a plan made several
moves earlier can all affect the time spent. None of them is contained in a
single position with its clocks.
\item \textbf{A point prediction without a distribution.} The model predicts
one value and not how widely players vary around it. A statement such as
``faster than nine in ten players of the same rating in this position''
requires the distribution, for example from a model that predicts quantiles.
\item \textbf{Whole seconds.} The clocks are recorded in whole seconds, so
the fastest moves are known only coarsely.
\item \textbf{Little data on time pressure.} Rapid games are played mostly
with ample time, and the groups with less than a minute on the clock are
small.
\item \textbf{No engine evaluation, by design.} The model does not know which
positions are critical. It measures how long players take, and not how long a
position deserves.
\item \textbf{Only analysed games.} The dataset is shared with a task that
needs engine evaluations, and Lichess analyses a game only when a user
requests it.
\item \textbf{One site, one time control, one month.} The games are the
first qualifying rapid games of one month on Lichess. A player's use of the
clock on another site or at another time control requires a separate
evaluation.
\item \textbf{Single training runs.} Each network was trained with one
random seed. The confidence intervals reflect the sampling of test games and
not the variation between training runs.
\end{itemize}

\section{Reproducibility}
\label{sec:cost}

Training the \V{runs.count} models on the leaderboard took \V{cost.seconds}
in total: \V{cost.seconds.networks} for the networks, each trained on a
single cloud GPU, and \V{cost.seconds.others} for the constant baseline and
GBDT on a CPU. The ablation models that are not on the leaderboard are not
included in this total. The best model has \V{best.parameters} parameters and
was trained for \V{best.epochs.run} epochs in \V{best.seconds} on an
\V{best.device}, at about \V{best.rows-a-second} positions per second, not
counting the pretraining of its encoder. Table~\ref{tab:runs} lists
the hardware and training time of every model, and
Table~\ref{tab:settings} the hyperparameters of the networks of the fixed
sequence.

All results are generated from committed files. The dataset is built from
its recipe by \texttt{lab data game-positions}, the models are trained by
\texttt{lab all think-time}, each training run writes a record to the model's
\texttt{results/} directory, and \texttt{lab paper think-time} generates
every number, table and figure of this paper from those records and compiles
it. A record names the build of the dataset it used (\V{data.fingerprint})
and the commit of the code that produced it (\V{code.commit}). Model weights
and test predictions are not stored in the repository. The model's card
states where they are kept and gives the hash of the best model's weights.

\section{Conclusion and Future Work}
\label{sec:conclusion}

The time a player spends on a move in a rapid game can be predicted from the
position, the ratings and the clocks with an error factor of
\V{best.mae.factor}, compared with \V{floor.mae.factor} for a constant
prediction. Most of this improvement comes from the context, in particular
the time remaining on the mover's clock, and from a small number of
hand-crafted features. A network that reads the board adds a further, smaller
improvement, which is largest in positions with many captures. The
predictions are accurate on average over their whole range, and they leave
most of the variation between individual moves unexplained. The following
experiments are the natural next steps.
\begin{itemize}
\item \textbf{The distribution, not only its centre.} A model that predicts
a low, a middle and a high quantile of the time would let the application say
how unusual the time spent on a move was.
\item \textbf{The game as a sequence.} The preceding moves and the players'
use of their clocks contain what a single position does not: how long the
opponent has just thought, and how the player has used the clock so far.
Allie reads a game in this way~\cite{allie}.
\item \textbf{A joint model of thinking time and blunders.} The blunder-risk
model trained on the same positions underestimates the risk of the moves that
take longest. A model trained on both targets tests whether a long think can
be anticipated from the position.
\ifmore
\item \textbf{More data together with more optimisation steps.} On the
larger training sets the transformer reached its best validation MAE in the
final epoch of its schedule. Training on four times the games with a longer
schedule, and training the pretrained model on four times the games, are the
two missing cases.
\fi
\ifshapes
\ifchangestart
\item \textbf{The accepted modifications in combination.} The attack-aware
attention bias and the ResNet encoder were evaluated from random
initialisation only. A pretrained transformer with the attack-aware bias, and
a pretrained ResNet, have not been trained.
\fi
\fi
\ifchangestart
\item \textbf{A stronger pretrained encoder.} The encoder was pretrained
once, at one size, on part of the evaluation database. Longer pretraining,
more positions and a larger encoder have not been tried.
\fi
\item \textbf{A player's own use of the clock.} Before the application
displays a reference time to a player, the predictions should be compared
with the times in that player's own games.
\end{itemize}

\bibliographystyle{unsrt}
\bibliography{../../../papers/references}

\clearpage
\appendix

\section{Hyperparameters and Features}

\begin{table}[htbp]\centering\footnotesize
\block{settings}
\caption{Hyperparameters of the networks of the fixed sequence, from their
configuration files. \emph{LR} is the peak learning rate. Values shared by
all of them are given below the table.}
\label{tab:settings}
\end{table}

\begin{table}[htbp]\centering\footnotesize
\block{hand}
\caption{Hand-crafted features: the inputs of \V{task.reference} in addition
to the scalar context features, and of a network whose configuration includes
them.}
\label{tab:hand}
\end{table}

\iflook
\section{Examples with Extreme Predictions}

The three positions with the highest predictions were all played with an
initial time of 15 minutes and an increment of 10 seconds, with both queens
and most of the pieces on the board, and in each of them the time spent was
longer than the prediction (Figure~\ref{fig:highest}). Two of the three
positions with the lowest predictions are endgames in which the mover has
less than ten seconds left, and in the third the move played is a recapture.
All three of these moves were played within one second
(Figure~\ref{fig:lowest}).

\examples{examples-highest}{The test positions with the longest predicted
times of the best model.}{fig:highest}

\examples{examples-lowest}{The test positions with the shortest predicted
times of the best model.}{fig:lowest}

\section{Subgroup Tables}

Each table below partitions the test positions by one variable. The columns
are the group, the number of test positions, the typical time spent in
seconds (\emph{Observed}), the mean prediction in seconds and the MAE of each
of the two models compared, and the paired gain of the first model over the
second with its 95\% CI, in bold when statistically significant.

\block{slices}
\fi

\end{document}
